% ============================================================
% Echo-Search — Relazione di progetto (v3)
% Compilare con: latexmk -pdf echo-search-relazione-v3.tex
% ============================================================
\documentclass[11pt,a4paper]{report}

\usepackage[utf8]{inputenc}
\usepackage[T1]{fontenc}
\usepackage[italian]{babel}
\usepackage{lmodern}
\usepackage{microtype}
\usepackage[margin=2.6cm]{geometry}
\usepackage{setspace}
\setstretch{1.18}
\usepackage{booktabs}
\usepackage{tabularx}
\usepackage{amsmath}
\usepackage{amssymb}
\usepackage{xcolor}
\usepackage{tikz}
\usetikzlibrary{arrows.meta, positioning, fit, backgrounds}
\usepackage{listings}
\usepackage{enumitem}
\setlist{itemsep=1pt, topsep=3pt, parsep=0pt}
\usepackage[hidelinks]{hyperref}

\definecolor{echoblue}{HTML}{2563EB}
\definecolor{echoamber}{HTML}{EA7317}
\definecolor{echoslate}{HTML}{64748B}
\definecolor{echogreen}{HTML}{16A34A}
\definecolor{echored}{HTML}{DC2626}
\definecolor{codebg}{HTML}{F4F6F9}

\lstset{basicstyle=\ttfamily\scriptsize, backgroundcolor=\color{codebg},
  breaklines=true, frame=single, rulecolor=\color{echoslate!40},
  columns=fullflexible, keepspaces=true, xleftmargin=3pt, xrightmargin=3pt,
  literate={à}{{\`a}}1 {è}{{\`e}}1 {é}{{\'e}}1 {ì}{{\`i}}1 {ò}{{\`o}}1 {ù}{{\`u}}1}

\tikzset{
  box/.style={draw=echoblue, fill=echoblue!7, rounded corners=2pt, align=center,
              minimum height=0.8cm, font=\small, thick},
  io/.style={draw=echoamber, fill=echoamber!12, rounded corners=2pt, align=center,
             minimum height=0.8cm, font=\small, thick},
  metric/.style={draw=echogreen, fill=echogreen!12, rounded corners=2pt, align=center,
                 minimum height=0.8cm, font=\small, thick},
  arr/.style={-{Latex[length=1.8mm]}, echoslate, thick},
  lbl/.style={font=\scriptsize\itshape, color=echoslate},
}

\newcommand{\avvertenza}[2]{%
  \begin{center}
  \fcolorbox{echored!50}{echored!5}{%
    \begin{minipage}{0.94\textwidth}
    \small\textbf{#1}\\[2pt] #2
    \end{minipage}}
  \end{center}}

\newcommand{\ES}{\textbf{Echo-Search}}

% Capitoli più compatti
\usepackage{titlesec}
\titleformat{\chapter}[hang]{\LARGE\bfseries}{\thechapter.}{12pt}{\LARGE}
\titlespacing*{\chapter}{0pt}{-18pt}{18pt}

\begin{document}

% ---------------- FRONTESPIZIO ----------------
\begin{titlepage}
\centering
\vspace*{1.5cm}
{\Huge\bfseries Echo-Search\par}
\vspace{0.5cm}
{\Large Un sistema anti--echo--chamber per la ricerca\\ su temi divisivi\par}
\vspace{1.6cm}
{\large Relazione di progetto\par}
\vspace{0.2cm}
{\normalsize Information Retrieval e Natural Language Processing\par}
\vspace{1.8cm}
\begin{tikzpicture}
  \node[box, minimum width=2.6cm] (a) {Media\\\footnotesize news};
  \node[metric, minimum width=2.4cm, right=1.3cm of a] (m) {Stance\\Diversity};
  \node[io, minimum width=2.6cm, right=1.3cm of m] (b) {Social};
  \draw[arr] (a) -- (m); \draw[arr] (b) -- (m);
\end{tikzpicture}
\vfill
{\large Davide Cinelli \quad\textbar\quad Giuseppe Cuiuri\par}
\vspace{0.3cm}
{\normalsize Università della Calabria --- A.A. 2025/2026\par}
\vspace{0.3cm}
{\large Luglio 2026\par}
\end{titlepage}

\tableofcontents

% =================================================================
\chapter{Introduzione}\label{cap:intro}

\section{Contesto e problema}
Il dibattito pubblico sui temi controversi si svolge sempre più in ambienti informativi che,
per effetto combinato della personalizzazione algoritmica e dell'omofilia sociale, espongono
l'utente prevalentemente a contenuti concordi con le sue opinioni pregresse --- il fenomeno
noto come \emph{echo chamber}, o \emph{filter bubble} nella sua variante
algoritmica~\cite{pariser,cinelli}. Su temi ad alto contenuto tecnico e forte carica
valoriale come l'energia nucleare, il cambiamento climatico e la vaccinazione, la formazione
di un'opinione informata richiederebbe accesso bilanciato alle argomentazioni di entrambe le
parti, corredate da evidenze verificabili.

Gli strumenti di ricerca, però, sono ottimizzati per la \emph{pertinenza} e non per la
\emph{diversità}. Dato un corpus $\mathcal{D}$, una query $q$ e un \emph{target} di dibattito
$t$, un motore classico restituisce
\begin{equation}
R_k(q) = \operatorname*{arg\,top-}k_{\,d \in \mathcal{D}} \; \text{rel}(q, d),
\end{equation}
e nulla in questa formulazione vincola la distribuzione delle posizioni
$s(d,t)$ all'interno di $R_k(q)$, dove
$s(d,t) \in \{\textsc{favor}, \textsc{neutral}, \textsc{against}\}$: se fra i documenti più
pertinenti ce ne sono 18 favorevoli e 2 contrari, il sistema ne mostra 18 e 2. Poiché la query stessa porta con sé un'inclinazione implicita
(``\emph{risks} of nuclear energy'' e ``\emph{benefits} of nuclear energy'' pescano da regioni
diverse del corpus), il ranking per pertinenza è un amplificatore involontario dell'echo
chamber.

Il problema ha due livelli, che questo lavoro tiene distinti. Al livello
dell'\textbf{esposizione} --- quali documenti vengono mostrati --- è possibile intervenire,
diversificando il result set. Al livello dell'\textbf{ambiente} --- quali documenti esistono
nel corpus, e con quale distribuzione di posizioni --- non c'è intervento possibile: c'è una
proprietà da misurare. \ES{} si colloca all'intersezione: sostituisce l'output piatto con tre
colonne di prospettiva contrapposta, e misura contestualmente quanto plurale sia il materiale
che il corpus rende disponibile per quella query.

\section{Domande di ricerca}\label{sec:rq}
\begin{description}[leftmargin=1.5cm,style=nextline]
  \item[RQ1 --- Misurabilità] È possibile operazionalizzare l'omogeneità di opinione di un
        insieme di risultati in una metrica scalare, interpretabile e confrontabile fra corpora?
  \item[RQ2 --- Confronto fra ecosistemi] L'ecosistema social e quello dei media differiscono,
        sugli stessi temi, per diversità di stance del materiale che rendono recuperabile? In
        quale direzione?
  \item[RQ3 --- Motore di verdetto] Nel fact-checking \emph{a parità di evidenze}, un LLM
        generalista di taglia contenuta ($\le$10B) che \emph{ragiona} è competitivo con un
        classificatore NLI specializzato e FEVER-tuned?
  \item[RQ4 --- Fattibilità] È possibile realizzare l'intera pipeline in modo locale, offline
        e riproducibile entro 6\,GB di memoria video?
\end{description}

Anticipiamo il grado di risposta raggiunto: RQ1, RQ3 e RQ4 ricevono risposta affermativa e
documentata; RQ2 riceve una risposta \emph{condizionata}, la cui validità interna è limitata
dal modo in cui i due corpora sono stati costruiti --- limite analizzato nella
Sezione~\ref{sec:validita-interna} e non nascosto.

\section{Contributi}
\begin{enumerate}[leftmargin=*]
  \item \textbf{Una pipeline simmetrica a due ecosistemi}: la medesima catena, con gli stessi
        modelli e parametri, su due corpora indicizzati separatamente. È la simmetria a rendere
        confrontabili i risultati.
  \item \textbf{Una metrica di echo chamber}, la \emph{Stance Diversity}, con l'analisi delle
        sue proprietà e delle condizioni sotto cui il confronto fra due valori è lecito.
  \item \textbf{Un fact-checker locale a due motori intercambiabili}: stessa knowledge base,
        stesso recupero, stesso cancello di pertinenza; cambia solo il modulo che emette il
        verdetto. È questo isolamento a rendere il confronto un esperimento controllato.
  \item \textbf{Uno stance detector fine-tunato}, valutato su benchmark citabili, in
        \emph{cross-dataset} e su un gold set nel dominio applicativo.
  \item \textbf{Una valutazione del retrieval su qrels propri}, doppiamente annotati da due
        persone in cieco, con ablazione a sei sistemi e misura di quanto l'ordinamento dipenda da
        chi ha giudicato.
\end{enumerate}

\section{Che cosa questo lavoro non è}
Non è un contributo modellistico: nessuna architettura nuova, nessuna funzione di perdita
nuova; l'unico addestramento è un fine-tuning supervisionato standard. Non è uno studio
sociologico sulla polarizzazione: i corpora sono campioni di dimensione limitata, scelti per
rendere il sistema eseguibile e riproducibile, non per essere rappresentativi dell'ecosistema
mediatico. Le conclusioni sulla differenza fra media e social vanno lette come proprietà
\emph{dei corpora costruiti}. Non è infine un sistema di moderazione: il registro di
credibilità delle fonti è usato per descrivere i dati, mai come filtro automatico sui
risultati mostrati.

\section{Struttura}
Il Capitolo~\ref{cap:sota} richiama il background necessario a giustificare le scelte. Il
Capitolo~\ref{cap:sistema} descrive architettura e implementazione. Il
Capitolo~\ref{cap:dati} documenta i dati --- capitolo cruciale, perché è lì che nascono sia i
risultati sia i loro limiti. Il Capitolo~\ref{cap:metodo} formalizza la metodologia; il
Capitolo~\ref{cap:protocollo} definisce il protocollo sperimentale \emph{prima} di vederne i
risultati, che sono nel Capitolo~\ref{cap:risultati}. Il Capitolo~\ref{cap:discussione}
discute findings, errori e minacce alla validità; il Capitolo~\ref{cap:conclusioni} conclude.

% =================================================================
\chapter{Background}\label{cap:sota}

\section{Echo chamber e diversità dei risultati}
Il termine \emph{echo chamber} descrive un ambiente in cui le convinzioni sono amplificate per
ripetizione, con esposizione limitata a punti di vista alternativi; la variante \emph{filter
bubble}~\cite{pariser} enfatizza la componente algoritmica. Sul piano empirico l'ambito
computazionale ha proceduto su due direttrici quasi sempre separate: la \emph{misura} della
polarizzazione --- tipicamente su grafi di interazione sociale, come in Cinelli et
al.~\cite{cinelli} --- e l'\emph{intervento}, cioè l'esposizione deliberata a contenuti
discordanti. Misurare la polarizzazione direttamente sui \emph{contenuti testuali}, tramite
stance detection, è una strada meno battuta e più direttamente collegabile a un motore di
ricerca.

La diversificazione dei risultati ha una tradizione consolidata nell'IR, nata però per un
problema diverso: l'ambiguità della query. In \emph{Maximal Marginal Relevance}~\cite{mmr} si
premia la rilevanza e si penalizza la ridondanza \emph{testuale}; approcci successivi
esplicitano i sotto-argomenti da coprire. La differenza rispetto a \ES{} è sostanziale: là la
diversità è dissimilarità nello spazio delle rappresentazioni, qui è definita \emph{a priori}
e semanticamente --- le tre stance. Ne consegue che la misura è direttamente interpretabile e
confrontabile fra corpora, ma anche che la sua qualità dipende da quella del classificatore di
stance, con un errore di misura di cui va reso conto.

\section{Information Retrieval}

\subsection{Retrieval sparso}
La famiglia \emph{Best Matching} culmina in BM25~\cite{bm25}, tuttora una baseline difficile
da battere:
\begin{equation}\label{eq:bm25}
\text{BM25}(D,Q) = \sum_{t \in Q} \text{IDF}(t) \cdot
\frac{f(t,D)\,(k_1 + 1)}{f(t,D) + k_1\left(1 - b + b\,\dfrac{|D|}{\text{avgdl}}\right)} .
\end{equation}
I due iperparametri hanno interpretazioni distinte e non intercambiabili. Il parametro $k_1$
governa la \emph{saturazione} della term-frequency: per $k_1 \to 0$ il contributo di un
termine diventa binario, per $k_1$ grande cresce quasi linearmente. Il parametro $b \in [0,1]$
governa la \emph{normalizzazione per lunghezza}: con $b=0$ la lunghezza è ignorata, con $b=1$
la normalizzazione è piena. Il default $b = 0{,}75$ è tarato su collezioni web eterogenee;
questa osservazione motiva la scelta adottata nella Sezione~\ref{sec:metodo-retrieval}. Il
limite strutturale di BM25 è il \emph{vocabulary mismatch}: essendo lessicale, non riconosce
che ``meltdown'' e ``nuclear accident'' denotano lo stesso evento --- limite tanto più
rilevante su un dominio argomentativo, dove la stessa posizione si esprime con lessici molto
diversi a seconda dell'orientamento della fonte.

\subsection{Retrieval denso e indici vettoriali}
I metodi \emph{dense} proiettano i testi in uno spazio semantico tramite \emph{bi-encoder}:
query e documento sono codificati \emph{separatamente} e la rilevanza è approssimata dalla
similarità coseno, con addestramento per \emph{contrastive learning} su grandi collezioni di
coppie~\cite{dpr}. Il vantaggio è decisivo: i documenti si codificano offline una volta sola.
Il modello adottato, \emph{bge-base-en-v1.5}~\cite{bge}, è \emph{asimmetrico}: prevede un
prefisso-istruzione sulla query e nessun prefisso sui passaggi. Ignorare l'asimmetria degrada
silenziosamente il recupero senza generare alcun errore --- difetto insidioso e frequente.

Su vettori L2-normalizzati il prodotto interno coincide con la similarità coseno, e librerie
come FAISS~\cite{faiss} offrono indici \emph{esatti} o \emph{approssimati}. Per collezioni
dell'ordine di $10^3$--$10^4$ vettori l'indice esatto è preferibile: il costo è trascurabile e
si ottiene in cambio la riproducibilità perfetta dei risultati, requisito non negoziabile in
un contesto di valutazione.

Il limite dei bi-encoder è l'assenza di interazione fra query e documento, che il secondo
stadio è chiamato a recuperare.

\subsection{Retrieval ibrido e reranking}
Poiché segnale lessicale e semantico sono complementari, lo stato dell'arte li combina, con due
strategie dalle proprietà diverse. La \emph{fusione di punteggi} normalizza e combina
linearmente,
\begin{equation}\label{eq:fusion}
\text{score}(d) = \alpha\,\widetilde{s}_{\text{sparse}}(d) + \beta\,\widetilde{s}_{\text{dense}}(d),
\qquad \alpha + \beta = 1,
\end{equation}
conservando l'informazione sulla \emph{magnitudo} delle differenze, al prezzo di dipendere
dalla normalizzazione e di richiedere la taratura di $\alpha$. La \emph{fusione di ranghi}, di
cui \emph{Reciprocal Rank Fusion}~\cite{rrf} è il rappresentante più usato,
\begin{equation}\label{eq:rrf}
\text{RRF}(d) = \sum_{r \in \mathcal{R}} \frac{1}{k_{\text{rrf}} + \text{rank}_r(d)},
\end{equation}
è robusta e priva di parametri, ma perde quell'informazione: due documenti in posizione 1 e 2
contribuiscono quasi allo stesso modo, che siano separati da un abisso o da un niente.

Un secondo stadio di \emph{reranking} affina l'ordinamento con un \emph{cross-encoder}, che
riceve query e documento \emph{insieme} e ne modella l'interazione con attenzione
completa~\cite{sbert}: sistematicamente più preciso del bi-encoder, ma con costo $O(N)$ per
query e non precomputabile. Ne discende lo schema \emph{retrieve-then-rerank}: uno stadio
economico e ad alto recall su tutto il corpus, seguito da uno costoso e preciso su poche
decine di candidati. I due stadi non sono alternativi: la fusione \emph{recupera}, il reranker
\emph{riordina}, e il tetto di prestazione del secondo è fissato dal recall del primo.

\subsection{Valutazione}
Il paradigma dominante resta quello di Cranfield: una collezione, un insieme di query e un
insieme di \emph{giudizi di rilevanza} (\emph{qrels}). Poiché giudicare l'intero corpus è
impraticabile, si ricorre al \emph{pooling}: si uniscono i primi $k$ risultati di tutti i
sistemi in esame e si giudica solo l'unione, assumendo non rilevanti i documenti mai entrati
nel pool. Le metriche standard sono il \emph{Normalized Discounted Cumulative Gain}
\begin{equation}\label{eq:ndcg}
\text{DCG@}k = \sum_{i=1}^{k} \frac{\text{rel}_i}{\log_2(i+1)}, \qquad
\text{NDCG@}k = \frac{\text{DCG@}k}{\text{IDCG@}k},
\end{equation}
che premia i rilevanti nelle prime posizioni ed è definito per giudizi graduati; la \emph{Mean
Average Precision}; e il \emph{Recall@k}, che è la metrica corretta per valutare uno stadio di
generazione dei candidati. L'accordo fra annotatori si riporta convenzionalmente con il
$\kappa$ di Cohen~\cite{cohen}.

\section{Stance detection}
La stance detection classifica la posizione di un testo verso un \emph{target} in FAVOR,
AGAINST, NEUTRAL. Va distinta dalla \emph{sentiment analysis}: un testo può esprimere
sentimento negativo ed essere a favore del target (``è \emph{tragico} che la disinformazione
stia rallentando le vaccinazioni''). I benchmark di riferimento sono \textbf{SemEval-2016
Task~6}~\cite{semeval}, su tweet e cinque target predefiniti, e \textbf{VAST}~\cite{vast}, con
migliaia di target diversi e quindi riferimento per la generalizzazione \emph{cross-target}.
L'approccio dominante fine-tuna transformer pre-addestrati su input a coppia
\texttt{[CLS] target [SEP] testo [SEP]}: la codifica a coppia è preferibile alla
concatenazione manuale perché sfrutta i \emph{token type embeddings} nativi.

Tre difficoltà sono documentate in letteratura e si ripresentano puntualmente in questo
lavoro: la \textbf{dipendenza dal target} (la stringa scelta per rappresentarlo influenza
sensibilmente le prestazioni); la classe \textbf{NEUTRAL}, notoriamente la peggio classificata
perché raccoglie fenomeni eterogenei --- testi informativi, ambigui, ironici, che riportano
posizioni altrui; e la \textbf{sensibilità al dominio}. Un'alternativa senza addestramento è
l'uso di modelli NLI in \emph{zero-shot}~\cite{yin}, adottata qui come baseline.

\section{Fact-checking e Natural Language Inference}
La verifica automatica di affermazioni è convenzionalmente scomposta in tre fasi, secondo
l'architettura di FEVER~\cite{fever}: estrazione della \emph{claim}, recupero delle
\emph{evidenze}, formulazione del \emph{verdetto}. La scomposizione ha un pregio metodologico:
rende ogni fase valutabile in isolamento, evitando che un errore di recupero venga attribuito
al motore di verdetto.

La terza fase è tipicamente modellata come \emph{Natural Language Inference}: data una coppia
(premessa, ipotesi) si predice se la prima \emph{implica}, \emph{contraddice} o è
\emph{neutrale} rispetto alla seconda, con corrispondenza diretta ai verdetti
SUPPORTED / REFUTED / NOT\_ENOUGH\_INFO. \textbf{ClimateFEVER}~\cite{climatefever} adatta il
paradigma a claim climatiche reali: benchmark sensibilmente più difficile di FEVER, perché le
claim sono formulate in linguaggio non controllato, spesso vaghe o parzialmente vere.

Più recentemente i \emph{Large Language Model} vengono impiegati come motore di verdetto,
sfruttandone la capacità inferenziale generale al posto della specializzazione. Il rischio noto
è che il modello risponda sulla base della conoscenza parametrica anziché delle evidenze
fornite, il che invaliderebbe il confronto con l'NLI e aprirebbe la porta alle allucinazioni;
la contromisura standard è il \emph{grounding} esplicito nel prompt. Il paradigma generale del
recupero di evidenze a supporto della generazione è quello della \emph{Retrieval-Augmented
Generation}~\cite{rag}.

\section{Misure di diversità}\label{sec:sota-diversity}
Per quantificare l'eterogeneità di una distribuzione categoriale lo strumento canonico è
l'\textbf{entropia di Shannon}~\cite{shannon}, $H(p) = -\sum_c p_c \log p_c$, massima e pari a
$\log K$ per la distribuzione uniforme e nulla per una distribuzione degenere. La
normalizzazione $H(p)/\log K$ produce un valore in $[0,1]$ interpretabile e confrontabile.
Alternative possibili sono l'indice di Gini-Simpson, che pesa meno le categorie rare, e i
numeri di Hill. L'entropia normalizzata è preferibile qui perché è la scelta più diffusa e
quindi confrontabile con la letteratura, ha un'interpretazione informazionale precisa, e
penalizza fortemente le distribuzioni degeneri --- che è esattamente il comportamento
desiderato quando l'obiettivo è segnalare un'echo chamber.

Va però sottolineato un punto che il Capitolo~\ref{cap:discussione} riprende: l'entropia è
funzione della sola distribuzione e non sa nulla di \emph{come} quella distribuzione sia stata
generata. Confrontare $H(p_1)$ e $H(p_2)$ è informativo solo se i due campioni provengono da
procedure comparabili. Questa condizione, apparentemente ovvia, è il punto delicato dell'intero
lavoro.

\section{Posizionamento}
\ES{} non propone nuovi modelli: il suo contributo è di \emph{sistema} e di \emph{metodo di
misura}. I due elementi di originalità sono l'operazionalizzazione dell'echo chamber come
metrica scalare comparabile fra ecosistemi, calcolata sui contenuti e non sulle relazioni
sociali; e un banco di prova per il fact-checking in cui il motore di verdetto è l'\emph{unica}
variabile che cambia, il che consente di attribuire la differenza di prestazione al motore e
non alla catena di recupero.

% =================================================================
\chapter{Il sistema}\label{cap:sistema}

\section{Visione d'insieme}
Il sistema è un motore di retrieval a due stadi seguito da moduli NLP che arricchiscono i
documenti recuperati (Figura~\ref{fig:arch}). Il principio guida è la \textbf{simmetria fra
ecosistemi}: la medesima pipeline, con gli stessi modelli e iperparametri, viene eseguita su
ciascuno dei due ecosistemi (media/news e social); è il \emph{confronto} fra i due blocchi a
costituire l'output di ricerca.

La simmetria non è un dettaglio estetico ma un requisito sperimentale: se i due ecosistemi
fossero trattati diversamente, ogni differenza osservata sarebbe attribuibile tanto
all'ecosistema quanto al trattamento. Vincolare la pipeline a essere identica elimina per
costruzione un'intera classe di confondimenti --- ne resta un'altra, di natura diversa,
discussa nella Sezione~\ref{sec:validita-interna}.

\begin{figure}[t]
\centering
\resizebox{\textwidth}{!}{%
\begin{tikzpicture}[node distance=0.42cm and 0.6cm]
  \node[io, minimum width=1.7cm] (corpus) {Corpus\\\footnotesize news \textbar\ social};
  \node[box, right=of corpus, minimum width=2.2cm] (retr)
    {Retrieval ibrido\\\footnotesize fusione $\to$ rerank};
  \node[io, right=of retr, minimum width=1.7cm] (docs) {Documenti\\\footnotesize recuperati};
  \node[box, above right=0.10cm and 0.8cm of docs, minimum width=2.1cm] (stance)
    {Stance\\\footnotesize DeBERTa-v3 FT};
  \node[box, below right=0.10cm and 0.8cm of docs, minimum width=2.1cm] (fact)
    {Fact-check\\\footnotesize NLI \textbar\ LLM};
  \node[box, right=of stance, minimum width=1.7cm] (bal) {Balance\\\footnotesize 3 colonne};
  \node[metric, right=of bal, minimum width=1.8cm] (div) {Stance\\Diversity};
  \node[box, below=0.5cm of div, minimum width=1.8cm] (summ) {Summary\\\footnotesize FLAN-T5};
  \node[io, minimum width=1.9cm, minimum height=2.4cm, anchor=north west] (out)
    at ([xshift=0.7cm]div.north east)
    {Dashboard\\\footnotesize 3 colonne\\\footnotesize news $\leftrightarrow$ social};
  \begin{pgfonlayer}{background}
    \node[draw=echoamber, dashed, thick, rounded corners=4pt,
          fit=(stance)(fact), inner sep=5pt] (par) {};
  \end{pgfonlayer}
  \node[lbl, below=1pt of par.south, color=echoamber, align=center]
    {rami indipendenti\\sui \emph{medesimi} documenti};
  \draw[arr] (corpus) -- (retr);
  \draw[arr] (retr) -- (docs);
  \draw[arr] (docs.east) -- (stance.west);
  \draw[arr] (docs.east) -- (fact.west);
  \draw[arr] (stance) -- (bal);
  \draw[arr] (bal) -- (div);
  \draw[arr] (bal.south east) -- (summ.west);
  \draw[arr] (div) -- (out.west|-div);
  \draw[arr] (summ.east) -- (out.west|-summ);
  \draw[arr] (fact) -| (out.south);
\end{tikzpicture}}
\caption{Architettura del sistema. \textbf{Stance} e \textbf{fact-check} operano entrambi sui
documenti recuperati, in modo reciprocamente indipendente; \textbf{balance} consuma le stance e
\textbf{summary} le colonne bilanciate. Il fact-check non alimenta il riassunto: arricchisce i
documenti mostrati.}
\label{fig:arch}
\end{figure}

\section{I moduli NLP sono un DAG, non una catena}\label{sec:modularita}
Una precisazione architetturale importante: i quattro moduli NLP \emph{non} formano una catena
in cui l'output di ciascuno è l'input del successivo. Analizzando le dipendenze dei dati,
\textbf{stance} e \textbf{fact-check} operano entrambi direttamente sui documenti recuperati e
in modo reciprocamente indipendente --- il verdetto di fact-checking non richiede la stance, né
viceversa --- e sono quindi concettualmente parallelizzabili; \textbf{balance} dipende dalle
stance, che usa per comporre le colonne e calcolare la diversità; \textbf{summary} dipende dalle
colonne bilanciate.

La distinzione ha una conseguenza pratica: il cammino critico è
$\text{retrieval} \to \text{stance} \to \text{balance} \to \text{summary}$, mentre il
fact-checking --- di gran lunga il modulo più lento con il motore LLM --- non vi appartiene e
potrebbe essere sovrapposto. Nell'implementazione attuale i due rami sono eseguiti in sequenza.
Si osservi inoltre che il fact-check è applicato \emph{dopo} il bilanciamento e non a tutti i
documenti recuperati: non è una dipendenza logica ma un'ottimizzazione, poiché verificare
documenti che verranno scartati costerebbe il doppio senza beneficio per l'utente.

\section{Corpus a due ecosistemi}
Per ciascun tema si mantengono due corpora fisicamente separati e indicizzati in modo
indipendente. La separazione è essenziale: con un indice unico il retrieval scegliere\-bbe fra le
due popolazioni in base alla sola pertinenza, e un tweet di venti parole competerebbe con un
articolo di duemila con punteggi non commensurabili --- BM25 pesa la lunghezza, il denso soffre
l'asimmetria di granularità --- così che la composizione del result set dipenderebbe da artefatti
di lunghezza più che dal contenuto. Con indici distinti ogni ecosistema è interrogato al suo
interno e i due result set sono confrontabili come campioni omogenei.

\section{Scelte di progettazione e alternative scartate}
La Tabella~\ref{tab:decisioni} raccoglie le decisioni rilevanti, ciascuna con l'alternativa
considerata e la ragione dello scarto.

\begin{table}[p]
\centering
\small
\begin{tabularx}{\textwidth}{p{2.9cm}p{3.0cm}X}
\toprule
\textbf{Decisione} & \textbf{Alternativa scartata} & \textbf{Motivazione} \\
\midrule
Retrieval ibrido sparso $+$ denso &
Solo denso, oppure solo BM25 &
Complementari: BM25 cattura i termini tecnici esatti (``PWR reactor'', ``CO$_2$ ppm''), il denso
sinonimi e parafrasi del dibattito argomentativo. Presi singolarmente lasciano scoperte classi
di query complementari (verificato empiricamente, \S\ref{sec:ris-retrieval}). \\
\addlinespace
Indice FAISS esatto &
Indice approssimato (HNSW, IVF) &
Su $10^3$--$10^4$ vettori il vantaggio di latenza è irrilevante, mentre l'esattezza garantisce
riproducibilità perfetta. \\
\addlinespace
Un indice per (tema $\times$ ecosistema) &
Indice unico globale &
L'utente interroga un tema specifico; indici separati evitano il rumore cross-tema e sono il
prerequisito del confronto fra ecosistemi. \\
\addlinespace
Bilanciamento deterministico delle colonne &
Loop agentico con LLM orchestratore &
Il riempimento ordinato è $O(n)$, deterministico e ispezionabile. Un loop agentico
introdurrebbe non-determinismo e costo senza garanzie aggiuntive. \\
\addlinespace
Full fine-tuning di DeBERTa-v3-base &
Mixture-of-Experts; DANN; LoRA &
MoE introduce non-determinismo nel routing, rendendo la Macro-F1 non riproducibile su test
fisso. DANN aggiunge complessità di addestramento non giustificata quando la generalizzazione si
ottiene a monte tramite VAST. Il modello è piccolo (184M) e sta nei 6\,GB in fp16. \\
\addlinespace
Diversità calcolata sul pool &
Diversità sulle colonne mostrate &
Le colonne sono bilanciate per costruzione: misurarne l'entropia restituirebbe una costante
travestita da misura. \\
\addlinespace
Fact-checking su KB locale &
API di ricerca web o servizi cloud &
Riproducibilità (la KB è congelata e versionata per hash), assenza di quote e costi,
funzionamento offline. Il prezzo è la copertura limitata, dichiarata fra i limiti. \\
\addlinespace
Due motori di verdetto sullo stesso stack &
Due sistemi di fact-checking distinti &
Condividendo estrazione, recupero e cancello di pertinenza, l'unica variabile che cambia è il
motore: la differenza misurata è attribuibile ad esso. \\
\addlinespace
Beam search deterministico &
Sampling con temperatura &
La riproducibilità richiede che lo stesso cluster produca sempre lo stesso riassunto. Con
\texttt{do\_sample=False} la temperatura è ininfluente: usarla sarebbe una contraddizione
tecnica. \\
\bottomrule
\end{tabularx}
\caption{Decisioni architetturali, alternative considerate e motivazioni.}
\label{tab:decisioni}
\end{table}

\section{Il vincolo hardware come vincolo di progetto}\label{sec:memoria}
L'intero sistema è progettato per una GPU portatile da \textbf{6\,GB} (NVIDIA RTX~4050 Laptop).
Far coesistere un bi-encoder, un cross-encoder, un classificatore di stance, un modello NLI, un
modello generativo e --- nel percorso LLM --- un modello da 7 miliardi di parametri richiede
misure esplicite:

\begin{itemize}[leftmargin=*]
  \item \textbf{Caricamento pigro}: nessun modello è istanziato nel costruttore; ogni componente
        pesante è una \emph{property} caricata al primo uso. Una ricerca senza fact-checking né
        riassunti non carica mai spaCy, l'NLI o FLAN-T5. Il beneficio collaterale è la
        testabilità della logica pura senza scaricare gigabyte di pesi.
  \item \textbf{Cache a voce singola}: cambiare tema \emph{sostituisce} l'applicazione caricata
        anziché affiancarne una seconda.
  \item \textbf{Allocatore a segmenti espandibili} e liberazione esplicita della cache CUDA fra
        i due ecosistemi, per contenere la frammentazione --- causa più frequente di
        esaurimento apparente quando si alternano modelli di dimensioni diverse.
  \item \textbf{Parametri ridotti nel solo percorso interattivo} (batch, claim per documento,
        beam): degradazione consapevole di qualità in cambio di latenza. Non è applicata al
        percorso di valutazione, quindi i numeri del Capitolo~\ref{cap:risultati} non sono
        ottenuti con la configurazione ridotta.
\end{itemize}

Il vincolo ha avuto un effetto in parte benefico sulla qualità metodologica: ha reso
impraticabile l'accumulo indiscriminato di modelli sempre più grandi e ha imposto di verificare,
per ogni componente, se un modello maggiore valesse il suo costo --- l'ablazione base~vs~large
sul motore NLI (\S\ref{sec:ris-ablazione}) nasce da questa necessità.

\section{Organizzazione del codice}
Il codice è organizzato per responsabilità: \texttt{preprocessing}, \texttt{retrieval}
(sparse, dense, hybrid, reranker, column\_balancer), \texttt{nlp} (stance, fact-checker nelle
sue varianti, summarizer), \texttt{evaluation} (metriche e caricamento dei qrels),
\texttt{gui}. Gli script di costruzione dei dati e di valutazione stanno in \texttt{scripts/},
separati dal codice di libreria perché hanno dipendenze più pesanti e non devono essere
importati dalla pipeline online.

Tutti gli iperparametri risiedono in \texttt{configs/config.yaml}, con la motivazione di
ciascuno annotata accanto al valore; due file ausiliari definiscono temi, target e query
canoniche (\texttt{topics.yaml}) e il registro di credibilità delle fonti
(\texttt{sources.yaml}). Vi risiedono anche le tre soglie del fact-checker --- il cancello di
pertinenza e le due soglie asimmetriche del verdetto --- che sono scelte sperimentali e come tali
devono essere visibili e modificabili senza toccare il codice.

Una convenzione che si è rivelata utile è la separazione sistematica fra \textbf{funzioni pure}
e codice di I/O negli script: nella costruzione del corpus social, per esempio, parsing,
campionamento e costruzione dei documenti sono funzioni pure, e solo \texttt{main()} accede a
rete e filesystem. Il beneficio è la testabilità: la suite conta \textbf{174 test} eseguiti in
meno di 30 secondi, senza rete e senza GPU, e copre preprocessing, costruzione dei dati,
aggregazione dei verdetti, parsing delle risposte LLM, rendering della GUI, metriche di ranking,
$\kappa$ e Stance Diversity.

\section{Interfaccia}
La dashboard (Streamlit) traduce le scelte metodologiche in scelte di presentazione: tre colonne
affiancate con palette divergente; una barra di diversità segmentata che mostra
contemporaneamente il punteggio e le proporzioni del pool; schede per ecosistema, così che il
confronto sia un cambio di scheda a parità di tutto il resto; card espandibili con claim,
verdetto e collegamenti alle fonti. Il \emph{soft-filter} merita una nota: selezionando una
posizione la colonna corrispondente si allarga, ma le altre \emph{non} scompaiono --- un filtro
rigido ricostruirebbe esattamente l'echo chamber che il sistema intende contrastare.

Le query canoniche sono \textbf{precalcolate}: corpus fisso e modelli deterministici implicano
che i risultati non cambino, quindi si calcolano una volta e si rileggono da JSON in
millisecondi, senza GPU a runtime. La motivazione è pratica: il fact-checking con motore LLM
richiede minuti per ecosistema, il che lo renderebbe inadatto a una dimostrazione interattiva.

% =================================================================
\chapter{Dati}\label{cap:dati}

Questo capitolo documenta i dati impiegati. Riceve spazio ampio perché è dove nascono
contemporaneamente i risultati del lavoro e i loro limiti: le decisioni di campionamento prese
qui determinano che cosa i numeri del Capitolo~\ref{cap:risultati} possano dimostrare.

\section{Temi e target}\label{sec:target}
Sono stati selezionati tre temi divisivi: energia nucleare, cambiamento climatico, vaccini. Per
ciascuno sono definiti parole chiave per il filtraggio, un \emph{target} per la stance detection
e un insieme di query canoniche (Tabella~\ref{tab:temi}).

La distinzione fra \emph{tema} e \emph{target} merita attenzione, perché è una delle poche
scelte tarate empiricamente. Il target è la stringa passata allo stance detector come primo
segmento della coppia, e la sua formulazione influenza il comportamento del classificatore. Due
decisioni sono state prese su misure e non su intuizioni. Per il clima si usa
\texttt{"global warming"} anziché \texttt{"climate change"}: il primo è più \emph{marcato}, dato
che nel dibattito anglosassone ``climate change'' è divenuto largamente descrittivo mentre
``global warming'' conserva la carica polemica delle posizioni negazioniste; la sostituzione ha
portato il recall su AGAINST dal $51\%$ al $78\%$. Per i vaccini si usa il target generale
\texttt{"vaccine safety"} anziché \texttt{"COVID-19 vaccination"}: quest'ultimo produceva una
diversità sul corpus news implausibilmente alta ($0{,}69$, superiore a quella del social) senza
alcun miglioramento del recall su AGAINST.

Il secondo caso è metodologicamente istruttivo: un target mal scelto può \emph{gonfiare} la
metrica di diversità senza che nulla sia migliorato. È un promemoria del fatto che la Stance
Diversity è una misura derivata, la cui affidabilità non supera quella del classificatore che la
alimenta.

\begin{table}[t]
\centering
\small
\begin{tabularx}{\textwidth}{lllX}
\toprule
\textbf{Tema} & \textbf{Target} & \textbf{Ecosistemi} & \textbf{Query canoniche} \\
\midrule
Nucleare & \texttt{nuclear energy} & news &
\emph{risks of nuclear energy}; \emph{is nuclear power safe}; \emph{nuclear energy and climate
change}; \emph{nuclear waste disposal} \\
\addlinespace
Clima & \texttt{global warming} & news, social &
\emph{is climate change caused by humans}; \emph{economic cost of climate policy};
\emph{climate change and extreme weather} \\
\addlinespace
Vaccini & \texttt{vaccine safety} & news, social &
\emph{are vaccines safe}; \emph{vaccine side effects}; \emph{vaccine mandates};
\emph{do vaccines cause autism} \\
\bottomrule
\end{tabularx}
\caption{Temi, target di stance, ecosistemi disponibili e query canoniche.}
\label{tab:temi}
\end{table}

L'ultima query dei vaccini è inclusa deliberatamente come banco di prova del fact-checker:
verte su una claim falsa notoria di cui la knowledge base contiene la confutazione esplicita.

\section{Corpus news: CC-News}
L'ecosistema media/news è costruito da \textbf{CC-News} (\texttt{vblagoje/cc\_news}), consumato
in streaming. Ogni articolo attraversa una catena di filtri: lunghezza minima (i frammenti
troncati non hanno contenuto argomentativo); un filtro anti-spam commerciale; l'assegnazione del
tema per \emph{densità} di parole chiave, dove un articolo è accettato solo se supera
contemporaneamente una soglia sul numero di occorrenze e una sulla loro densità --- la seconda è
essenziale, perché un articolo di tremila parole che cita ``climate change'' una volta di
passaggio non è un articolo sul clima; e infine un \textbf{tetto del 15\% per dominio}, senza il
quale le agenzie ad alto volume saturerebbero il corpus, con il duplice danno di ridurne
artificialmente la diversità e di sovrarappresentare articoli quasi duplicati.

\begin{table}[t]
\centering
\small
\begin{tabularx}{\textwidth}{lrrrX}
\toprule
\textbf{Tema} & \textbf{Doc.} & \textbf{Tier 2} & \textbf{Tier 3} & \textbf{Domini più frequenti} \\
\midrule
Nucleare & 403 & 38 & 365 & uk.reuters.com (46), reuters.com (19), dailymail.co.uk (13) \\
Clima & 800 & 49 & 751 & uk.reuters.com (79), eurekalert.org (30), channelnewsasia.com (28) \\
Vaccini & 714 & 53 & 661 & nbcnews.com (19), dailymail.co.uk (15), uk.reuters.com (13) \\
\midrule
\textbf{Totale} & \textbf{1917} & \textbf{140} & \textbf{1777} & \\
\bottomrule
\end{tabularx}
\caption{Composizione del corpus news dopo il filtraggio (708\,241 articoli scansionati).}
\label{tab:corpus-news}
\end{table}

Due osservazioni sulla Tabella~\ref{tab:corpus-news}. Il tema nucleare non raggiunge la quota di
800 documenti: la finestra temporale di CC-News non contiene abbastanza articoli che superino
congiuntamente le soglie di densità e il tetto per dominio. La colonna \emph{Tier} indica invece
che la maggioranza dei documenti ricade nel tier~3, che è il valore \emph{prudenziale di default}
assegnato ai domini non presenti nel registro delle fonti: non significa che il corpus sia di
bassa qualità, ma che il registro copre una frazione minoritaria dei domini presenti.

\paragraph{Un dato rilevante per l'interpretazione.}\label{sec:against-news}
Una misura condotta sul corpus news con il modello di stance fine-tunato rileva che i documenti
classificati AGAINST costituiscono il \textbf{$2{,}6\%$} del corpus clima e il
\textbf{$2{,}9\%$} del corpus vaccini. In positivo, è una misura diretta e non confondata
dell'omogeneità del corpus news: la cronaca su questi temi, nel campione considerato, quasi non
contiene posizioni contrarie. In negativo, è precisamente ciò che ha reso necessaria la
costruzione di un corpus social con campionamento stratificato --- e che ha così generato il
confondimento discusso nella sezione seguente.

\section{Corpus social}\label{sec:corpus-social}
L'ecosistema social utilizza due dataset pubblici di tweet etichettati per stance, scaricati
direttamente da repository GitHub senza alcun ricorso alle API di Twitter (che oggi non
sarebbero né gratuite né riproducibili):

\begin{description}[leftmargin=1.1cm,style=nextline]
  \item[Clima] \small
        \texttt{edwardcqian/climate\_change\_sentiment}, file \texttt{data/sample\_data.csv}.\\
        Campo testuale \texttt{message}, etichetta \texttt{sentiment}, mappatura
        $\{1 \to \textsc{favor},\, -1 \to \textsc{against},\, 0 \to \textsc{neutral}\}$; la
        classe ``2'' (segnalazione di notizia) è scartata perché non esprime una posizione.
  \item[Vaccini] \small
        \texttt{sohampoddar26/covid-vax-stance}, file
        \texttt{dataset/our\_labelled\_data.csv}.\\
        Campo testuale \texttt{tweet}, etichetta \texttt{majority\_label}, mappatura
        $\{\text{ProVax} \to \textsc{favor},\, \text{AntiVax} \to \textsc{against},\,
        \text{Neutral} \to \textsc{neutral}\}$.
\end{description}

Gli URL esatti, i conteggi e l'hash SHA-256 dei file prodotti sono in
\texttt{data/corpora/}\allowbreak\texttt{MANIFEST\_social.json}. Il trattamento comprende
deduplicazione per testo normalizzato (necessaria per eliminare i retweet, che gonfierebbero le classi più virali) e
campionamento \textbf{stratificato per classe} con tetto uniforme e seed fissato: 999 tweet per
il clima (333 per classe) e 870 per i vaccini (290 per classe).

\avvertenza{Avvertenza metodologica --- confondimento fra ecosistema e campionamento.}{%
Il corpus social ha, \emph{per costruzione}, entropia di stance massima ($1{,}000$): le tre
classi sono campionate in parti esattamente uguali. Il corpus news ha invece la distribuzione
\emph{naturale} risultante dal filtraggio, fortemente sbilanciata (circa il $3\%$ di AGAINST).
Ne consegue che il confronto diretto fra la diversità misurata sui due ecosistemi \textbf{non
isola una proprietà degli ecosistemi}: una parte non quantificabile a priori della differenza
osservata è attribuibile alla differenza fra un campionamento stratificato e uno naturale.}

La stratificazione non è stata arbitraria: senza di essa non sarebbe stato possibile disporre di
un gold set bilanciato su cui valutare lo stance detector nella classe AGAINST --- valutazione
che si è rivelata uno dei risultati diagnostici più informativi del progetto. In altri termini,
la stratificazione era la scelta giusta per l'obiettivo di \emph{valutare lo stance detector} ed
è quella sbagliata per l'obiettivo di \emph{confrontare gli ecosistemi}; il corpus è stato
costruito una volta e usato per entrambi gli scopi, e questo è l'errore di disegno. La
contromisura proposta è nella Sezione~\ref{sec:lift}.

\section{Knowledge base, dataset di stance e benchmark}
La fonte di evidenza del fact-checker è una \textbf{knowledge base Wikipedia} di 1589 articoli,
costruita espandendo categorie seme pertinenti ai tre temi --- incluse deliberatamente
\emph{Climate change denial} e \emph{Vaccine hesitancy}, perché la KB deve contenere le voci che
descrivono e confutano le posizioni negazioniste, altrimenti il fact-checker non troverebbe
evidenze proprio per le claim che più contano. Di ciascun articolo si scarica
l'\textbf{introduzione} anziché il testo completo: per convenzione redazionale è un riassunto dei
fatti salienti, quindi la porzione a più alta densità di evidenza verificabile. L'accesso avviene
via API MediaWiki con \emph{batching} e User-Agent conforme; ogni evidenza conserva titolo e URL
canonico, il che rende ogni verdetto tracciabile dall'interfaccia.

Il modello di stance è addestrato sull'unione di \textbf{VAST}~\cite{vast} e
\textbf{SemEval-2016 Task~6}~\cite{semeval}, normalizzati in uno schema comune
(Tabella~\ref{tab:stance-data}). Due caratteristiche vanno sottolineate: l'elevatissimo numero di
\textbf{target distinti} in addestramento (4\,644, interamente dovuto a VAST), che è la proprietà
da cui deriva la capacità di generalizzare a temi non visti; e lo \textbf{sbilanciamento di
NEUTRAL} ($20{,}1\%$ contro il $42{,}5\%$ di AGAINST), che motiva l'uso della Macro-F1 come
metrica di selezione del modello.

\begin{table}[t]
\centering
\small
\begin{tabular}{lrrrrrrr}
\toprule
\textbf{Split} & \textbf{Esempi} & \textbf{FAVOR} & \textbf{AGAINST} & \textbf{NEUTRAL} &
\textbf{VAST} & \textbf{SemEval} & \textbf{Target} \\
\midrule
Train & 16\,097 & 6\,005 & 6\,849 & 3\,243 & 13\,477 & 2\,620 & 4\,644 \\
Test  & 4\,255 & 1\,248 & 1\,733 & 1\,274 & 3\,006 & 1\,249 & 763 \\
\bottomrule
\end{tabular}
\caption{Composizione degli split per la stance detection.}
\label{tab:stance-data}
\end{table}

La valutazione del fact-checking impiega \textbf{ClimateFEVER}~\cite{climatefever}. Ogni claim è
corredata di più evidenze con etichetta propria, oltre all'etichetta complessiva: la doppia
annotazione consente la valutazione a due granularità descritta nella
Sezione~\ref{sec:proto-fc}. Le claim DISPUTED sono escluse dalla valutazione a livello di claim,
poiché il sistema non prevede una classe corrispondente e includerle significherebbe conteggiare
come errori casi in cui nessuna risposta disponibile è corretta.

\section{Riproducibilità dei dati}
Ogni fase di costruzione produce un \texttt{MANIFEST.json} che registra sorgente esatta,
parametri, conteggi, distribuzione per classe e un \textbf{hash SHA-256} del file prodotto. I
file di dati non sono versionati --- sono rigenerabili --- mentre i manifest lo sono: chiunque
rieseguga uno script può verificare l'identità bit a bit del proprio risultato con quello usato
qui. Tutti i campionamenti usano un seed esplicito.

% =================================================================
\chapter{Metodologia}\label{cap:metodo}

\section{Preprocessing}
Il preprocessing serve esclusivamente al ramo sparso: il ramo denso opera sul testo grezzo,
perché i modelli di embedding sono addestrati su testo naturale. Tre scelte lo caratterizzano.
La \textbf{lemmatizzazione} anziché lo stemming, che produce forme troncate spesso non parole,
scelta anche perché le frasi del corpus vengono mostrate all'utente come anteprime ed evidenze.
La \textbf{preservazione delle negazioni} (\texttt{not}, \texttt{never}, \texttt{n't} e una lista
estesa), benché figurino nelle liste standard di stopword: ``Nuclear power is \emph{not} safe'' e
``Nuclear power is safe'' hanno stance opposte e token quasi identici una volta rimosse. E la
\textbf{simmetria} fra documenti in indicizzazione e query in interrogazione: un'asimmetria è
sorgente classica di \emph{vocabulary mismatch} artificiale.

\section{Retrieval, stadio 1: fusione}\label{sec:metodo-retrieval}

\paragraph{Ramo sparso.} BM25 (Equazione~\ref{eq:bm25}) con $k_1 = 1{,}2$, valore di default
standard adottato da Lucene ed Elasticsearch, e $b = 0{,}60$, deliberatamente inferiore al
default $0{,}75$. La ragione è di dominio: la normalizzazione per lunghezza esiste per
contrastare un effetto spurio --- un documento lungo contiene più termini e ha più probabilità di
contenere quelli della query per caso --- ma su una collezione omogenea di articoli
giornalistici, dove quasi tutti i documenti sono lunghi, l'effetto spurio è minore mentre il
costo resta: penalizzare i documenti lunghi significa penalizzare gli articoli più ricchi di
evidenze fattuali, che sono quelli di cui il fact-checking a valle ha bisogno.

\paragraph{Ramo denso.} I documenti sono segmentati in \emph{chunk} da 512 token con overlap di
128 ($25\%$). Il valore 512 corrisponde al limite del contesto degli encoder BERT-based: un testo
più lungo verrebbe silenziosamente troncato dall'encoder. L'overlap del $25\%$ --- coerente con
la pratica di Lewis et al.~\cite{rag} --- evita che un'affermazione a cavallo di due chunk vada
perduta da entrambi. L'implementazione opera a livello di parola con un fattore di conversione
stimato, approssimazione dichiarata. Ogni chunk è codificato con \emph{bge-base-en-v1.5} (768
dim.), L2-normalizzato e inserito in un indice FAISS \texttt{IndexFlatIP}; la query è codificata
con il prefisso-istruzione, i passaggi mai.

Poiché l'unità indicizzata è il chunk ma quella restituita è il documento, serve
un'aggregazione. Si adotta il \textbf{max-pooling}:
\begin{equation}\label{eq:maxpool}
\text{dense}(d) = \max_{c \,\in\, \text{chunk}(d)} \text{sim}(q, c).
\end{equation}
La motivazione è che un documento è rilevante se \emph{almeno una} sua porzione lo è: un articolo
che dedica un paragrafo esattamente al tema della query è rilevante anche se i restanti nove
decimi parlano d'altro. Una media penalizzerebbe sistematicamente i documenti lunghi,
reintroducendo dal lato denso la distorsione che dal lato sparso si è attenuata abbassando $b$.

\paragraph{Normalizzazione e fusione.} Entrambi i rami producono un punteggio per \emph{ogni}
documento. I due vettori sono normalizzati min-max in $[0,1]$ e combinati secondo
l'Equazione~\ref{eq:fusion} con $\alpha = 0{,}45$ e $\beta = 0{,}55$. La normalizzazione è
indispensabile perché BM25 non è limitato superiormente mentre il coseno vive in $[-1,1]$:
sommarli grezzi equivarrebbe a ignorare il ramo denso. La scelta $\alpha < \beta$ riflette la
natura semantica del dominio, mantenendo però un peso sostanziale al segnale lessicale, perché la
terminologia tecnica specifica è proprio ciò che i modelli densi tendono a generalizzare
eccessivamente. La selezione dei primi $k$ avviene \emph{una sola volta} sui punteggi fusi,
ottenendo un pool di 50 documenti; è implementata anche la baseline RRF
(Equazione~\ref{eq:rrf}). Se l'indice denso non è disponibile il sistema degrada a solo-BM25 con
un avviso, anziché fallire.

\section{Retrieval, stadio 2: reranking}
Il pool viene riordinato da un cross-encoder (\emph{bge-reranker-base}) che ripunteggia i primi
30 candidati: il flusso è un imbuto $|\mathcal{D}| \to 50 \to 30$. Due precisazioni. Il reranker
è \textbf{agnostico alla stance}: riordina per pura pertinenza, e il bilanciamento avviene
interamente dopo. Mescolare i due criteri renderebbe impossibile distinguere, nel risultato
finale, ciò che è stato scelto per pertinenza da ciò che è stato scelto per bilanciamento, e
soprattutto renderebbe circolare la misura di diversità. Inoltre il reranker non può recuperare
ciò che lo stadio 1 non ha proposto: il suo tetto è il recall dello stadio 1, ed è questa la
ragione per cui il Recall@$k$ è la metrica corretta per valutare il primo stadio e l'NDCG@10
quella per la pipeline completa.

\section{Stance detection}\label{sec:metodo-stance}
Dato un target $t$ e un documento $d$ il modulo stima $P(s \mid t, d)$ su tre classi,
restituendo la classe di massima probabilità, la sua confidenza e la distribuzione completa ---
conservata e propagata fino all'interfaccia, così che l'utente possa distinguere una
classificazione netta da una incerta.

Il modello è un \textbf{DeBERTa-v3-base}~\cite{deberta} (184M) sottoposto a \emph{full
fine-tuning} su VAST~$\cup$~SemEval, con input codificato come coppia nativa
\texttt{[CLS] target [SEP] documento [SEP]}. Gli iperparametri sono nella
Tabella~\ref{tab:iperparam-training}; due meritano un commento. La selezione del checkpoint su
\textbf{Macro-F1 di validazione} anziché sulla loss: con classi sbilanciate la loss è dominata
dalle maggioritarie, e selezionare su di essa produrrebbe modelli che sacrificano NEUTRAL. E la
combinazione \emph{batch size} 8 con \emph{gradient accumulation} 4, che realizza un batch
efficace di 32 restando entro i 6\,GB: la traduzione diretta del vincolo hardware in una scelta
di addestramento.

\begin{table}[t]
\centering
\small
\begin{tabular}{ll@{\qquad}ll}
\toprule
\textbf{Parametro} & \textbf{Valore} & \textbf{Parametro} & \textbf{Valore} \\
\midrule
Modello base & \texttt{deberta-v3-base} & Warmup ratio & 0,06 \\
Lunghezza massima & 256 token & Weight decay & 0,01 \\
Epoche & 3 (early stopping, patience 1) & Precisione & fp16 \\
Batch size & 8 ($\times$4 accumulation $=$ 32) & Metrica di selezione & Macro-F1 val. \\
Learning rate & $2 \times 10^{-5}$ & Validazione & 10\%, stratificata \\
\bottomrule
\end{tabular}
\caption{Iperparametri di fine-tuning dello stance detector (seed 42).}
\label{tab:iperparam-training}
\end{table}

In inferenza i documenti sono troncati alle prime 400 parole ($\approx$512 token). Il troncamento
privilegia l'inizio dell'articolo, scelta ragionevole nel dominio giornalistico dove la struttura
a piramide rovesciata colloca in apertura la tesi principale, ma va dichiarata: un articolo che
esprime la propria posizione solo in chiusura verrebbe classificato sulla sua parte descrittiva.
Se il modello fine-tunato non è presente su disco il sistema degrada a una pipeline
\emph{zero-shot} su modello NLI, che è anche la \textbf{baseline} contro cui il fine-tuning è
misurato.

\section{Bilanciamento e Stance Diversity}\label{sec:metodo-diversity}
I documenti rerankati sono classificati e distribuiti nelle tre colonne con un algoritmo
deliberatamente elementare e deterministico:

\begin{lstlisting}[language=Python]
classifica tutti i documenti del pool (un solo batch)
per ogni documento d, in ordine di pertinenza decrescente:
    c = stance(d)
    aggiungi d a full[c]                  # tutto il pool: serve alla diversity
    se len(colonne[c]) < per_column:      # per_column = 5
        aggiungi d a colonne[c]           # le colonne prendono i piu' pertinenti
diversity = entropia_normalizzata(full)
\end{lstlisting}

Poiché il pool è già ordinato per pertinenza, il primo documento di una data posizione è anche il
più pertinente di quella posizione: le colonne raccolgono i migliori rappresentanti di ciascuna
fazione senza ordinamenti aggiuntivi. Se una posizione è rara l'algoritmo la pesca scorrendo più
in profondità; se è assente la colonna resta vuota, e questa è un'informazione, non un difetto.
È la versione deterministica dell'idea di \emph{self-correction retrieval}: nessun loop agentico,
solo una selezione ordinata $O(n)$ ispezionabile.

La diversità è l'entropia di Shannon normalizzata sulla distribuzione delle stance:
\begin{equation}\label{eq:diversity}
\text{Diversity}(p) = -\frac{1}{\log 3}\sum_{c \in \{\textsc{f},\textsc{n},\textsc{a}\}}
p_c \log p_c \;\in\; [0,1],
\end{equation}
con la convenzione $0\log 0 = 0$.

\paragraph{Perché sul pool e non sulle colonne.} La diversità è calcolata sulla distribuzione
dell'\textbf{intero pool classificato}, non sulle tre colonne mostrate. Le colonne sono
bilanciate \emph{per costruzione}: se il pool contiene almeno cinque documenti per posizione, le
colonne conterranno $5{+}5{+}5$ e la loro entropia sarà esattamente $1$, indipendentemente da
qualunque proprietà del corpus. Misurare lì restituirebbe una costante travestita da misura. Ne
segue la lettura dei due numeri che l'interfaccia mostra affiancati: la barra rappresenta
\emph{ciò che il corpus offre}, le colonne \emph{ciò che il sistema mostra}, e la distanza fra i
due è la misura di quanto il bilanciamento sia intervenuto.

\paragraph{Proprietà e limiti.} La metrica vale $0$ se e solo se tutti i documenti condividono la
stessa posizione e $1$ se e solo se le tre sono equamente rappresentate
(Tabella~\ref{tab:diversity-esempi}). È \emph{invariante per permutazione} delle classi: un pool
dominato da NEUTRAL riceve lo stesso punteggio di uno dominato da FAVOR, benché il primo descriva
una copertura descrittiva e il secondo una schierata --- per questo l'interfaccia mostra sempre
anche la distribuzione. È \emph{insensibile alla dimensione} del campione, desiderabile per la
confrontabilità ma con la conseguenza che la metrica non segnala la propria incertezza: su pool di
30 documenti la varianza campionaria non è trascurabile. Ed è calcolata su etichette
\emph{predette}: un classificatore che confonde le classi tende a produrre distribuzioni più
uniformi di quelle reali, gonfiando la diversità. È la minaccia principale
(\S\ref{sec:validita-interna}).

\begin{table}[t]
\centering
\small
\begin{tabular}{lccl}
\toprule
\textbf{Distribuzione (F, N, A)} & \textbf{Proporzioni} & \textbf{Diversity} & \textbf{Lettura} \\
\midrule
$(30,0,0)$ & $(1{,}00;\,0;\,0)$ & $0{,}00$ & echo chamber perfetta \\
$(28,1,1)$ & $(0{,}93;\,0{,}03;\,0{,}03)$ & $0{,}27$ & copertura quasi unanime \\
$(20,8,2)$ & $(0{,}67;\,0{,}27;\,0{,}07)$ & $0{,}73$ & dominante, minoranza presente \\
$(15,10,5)$ & $(0{,}50;\,0{,}33;\,0{,}17)$ & $0{,}92$ & dibattito sbilanciato ma reale \\
$(10,10,10)$ & $(0{,}33;\,0{,}33;\,0{,}33)$ & $1{,}00$ & bilanciamento perfetto \\
\bottomrule
\end{tabular}
\caption{Comportamento della Stance Diversity su un pool di 30 documenti.}
\label{tab:diversity-esempi}
\end{table}

\paragraph{Una metrica correttiva: il diversity lift.}\label{sec:lift}
Come anticipato, il confronto diretto fra le diversità di due corpora campionati diversamente non
isola una proprietà degli ecosistemi. Sia $p^{\text{pool}}_q$ la distribuzione nel pool per la
query $q$ e $p^{\text{corpus}}$ quella dell'intero corpus; si definisce
\begin{equation}\label{eq:lift}
\text{Lift}(q) = \frac{H\bigl(p^{\text{pool}}_q\bigr)}{H\bigl(p^{\text{corpus}}\bigr)} .
\end{equation}
Il lift risponde a una domanda migliore: \emph{il retrieval amplifica o attenua lo sbilanciamento
già presente nel corpus?} Un valore $>1$ indica che il sistema recupera un insieme più plurale del
corpus medio; un valore $<1$ che il ranking per pertinenza sta concentrando le prospettive, cioè
che il motore stesso contribuisce all'echo chamber. Il pregio decisivo è di essere
\textbf{invariante rispetto alla distribuzione a priori}: due corpora campionati diversamente
producono lift confrontabili. Il calcolo richiede una passata batch dello stance detector
sull'intero corpus; l'implementazione è indicata fra gli sviluppi immediati.

\section{Fact-checking}\label{sec:metodo-factcheck}
Il fact-checker realizza l'architettura a tre fasi di FEVER con una scelta distintiva: le prime
due fasi sono condivise, la terza è \emph{intercambiabile}.

\paragraph{Fase 1 --- estrazione delle claim.} Dal testo si estraggono le frasi dichiarative con
predicato verbale tramite spaCy. Non tutte meritano verifica: si adotta un'euristica di
\emph{check-worthiness} ispirata a ClaimBuster~\cite{claimbuster} che pesa le frasi su entità
nominate e valori numerici, con peso maggiore alle entità --- una claim verificabile cita
tipicamente fatti concreti e attribuiti (``\emph{France} generates \emph{70\%} of its electricity
from nuclear power''), mentre una percentuale isolata senza soggetto è ambigua. In esercizio si
verifica la claim più check-worthy di ciascun documento mostrato, limitazione dettata dal costo.

\paragraph{Fase 2 --- evidenze e cancello di pertinenza.} Le frasi della KB sono codificate con
\emph{bge} e la ricerca avviene per similarità coseno. Il componente critico è il
\textbf{cancello di pertinenza}: se la similarità della migliore evidenza è inferiore a $0{,}62$
il verdetto è immediatamente NOT\_ENOUGH\_EVIDENCE e il motore non viene nemmeno invocato.
L'accorgimento è meno banale di quanto sembri e nasce da un'osservazione empirica: un recupero
semantico su una KB tematica restituisce \emph{sempre} qualcosa, e passare al motore un'evidenza
solo tematicamente vicina produce un'alta percentuale di falsi REFUTED, perché il modello NLI
legge come contraddizione una premessa che parla d'altro pur nello stesso ambito. Il cancello
trasforma questo fallimento silenzioso in un onesto ``non ho evidenze sufficienti''.

\paragraph{Fase 3a --- motore NLI.} Un classificatore FEVER-tuned
(\texttt{DeBERTa-v3-base-mnli-fever-anli}) predice per ciascuna coppia (evidenza, claim) le
probabilità di entailment, neutralità e contraddizione. La regola di decisione è deliberatamente
\emph{asimmetrica}:
\begin{equation}
\text{verdetto} =
\begin{cases}
\textsc{supported} & \text{se } \max_i P_{\text{ent}}^{(i)} \ge 0{,}75 \;\wedge\;
                                \max_i P_{\text{ent}}^{(i)} \ge \max_i P_{\text{con}}^{(i)},\\[3pt]
\textsc{refuted}   & \text{se } \max_i P_{\text{con}}^{(i)} \ge 0{,}95,\\[3pt]
\textsc{nei}       & \text{altrimenti.}
\end{cases}
\end{equation}
L'asimmetria fra le due soglie riflette un giudizio sui costi d'errore: dichiarare falsamente
``confutata'' un'affermazione vera è, in un sistema che si presenta come strumento di verifica,
molto più grave del dichiarare ``evidenza insufficiente'' --- il primo produce disinformazione con
l'autorità del sistema, il secondo solo un'astensione. Gli indici delle classi sono letti
dinamicamente da \texttt{model.config.id2label}, non hardcodati: un errore silenzioso classico se
si sostituisce il modello NLI con uno che ordina le classi diversamente.

\paragraph{Fase 3b --- motore LLM.} Un Large Language Model locale (\texttt{qwen2.5:7b} via
Ollama) ragiona sull'insieme delle evidenze e restituisce un verdetto strutturato in JSON con
motivazione (prompt in Appendice~\ref{app:prompt}). Tre aspetti sono rilevanti. Il
\textbf{grounding puro}: il prompt istruisce il modello a giudicare esclusivamente sulle evidenze
fornite, mai sulla conoscenza parametrica. Il vincolo riduce le allucinazioni e --- soprattutto
--- rende il confronto con l'NLI onesto: senza grounding, un'eventuale superiorità dell'LLM
sarebbe attribuibile alla conoscenza pregressa anziché alla capacità di ragionare sull'evidenza,
rendendo il confronto privo di significato. Il \textbf{determinismo}: temperatura $0$ e output
vincolato a JSON. E la \textbf{robustezza del parsing}: funzioni pure e testabili senza rete che
tentano il JSON, in caso di fallimento cercano un'etichetta riconoscibile nel testo, normalizzano
un'ampia tabella di sinonimi che i modelli producono spontaneamente, e in assenza di segnale
ripiegano su NOT\_ENOUGH\_EVIDENCE. Anche l'irraggiungibilità del server produce NEI anziché
un'eccezione: il fallimento è \emph{sicuro} per costruzione.

\section{Summarization}
La sintesi per prospettiva impiega FLAN-T5-base con strategia \emph{extract-then-abstract}: una
fase estrattiva punteggia le frasi della colonna con un BM25 a livello di frase e ne seleziona le
più salienti, riducendo l'input sotto la finestra di contesto e insieme la ridondanza fra
documenti che riportano gli stessi fatti. La fase astrattiva usa \textbf{beam search puro}
(\texttt{do\_sample=False}, \texttt{no\_repeat\_ngram\_size} $=3$): lo stesso cluster deve
produrre sempre lo stesso riassunto. Si segnala una precisazione tecnica fonte di confusione
frequente: la temperatura ha effetto \emph{solo} quando il campionamento è attivo, e associarla al
beam search è una combinazione incoerente, non una regolazione della creatività.

Una nota di onestà implementativa: era stata sperimentata l'iniezione di marcatori di riferimento
(\texttt{[Doc N]}) per rendere ogni affermazione tracciabile, e la tecnica è stata
\emph{abbandonata} perché FLAN-T5-base riproduce i marcatori come stringhe senza usarli
correttamente, producendo attribuzioni errate --- peggio dell'assenza di attribuzione. La
tracciabilità è garantita a livello di colonna. Il modulo non è stato fine-tunato per l'assenza di
riassunti di riferimento di dominio e perché la strategia estrattiva riduce il compito generativo
alla sola riscrittura.

% =================================================================
\chapter{Protocollo sperimentale}\label{cap:protocollo}

Questo capitolo definisce che cosa viene misurato, come e su quali dati, \emph{prima} di
riportarne i risultati. La separazione non è formale: serve a rendere verificabile che le
metriche non siano state scelte dopo aver visto i numeri.

\section{Principi e ambiente}
Il protocollo segue quattro principi: \textbf{benchmark citabili} dove possibile (VAST,
SemEval-2016 T6, ClimateFEVER) anziché dati costruiti ad hoc; \textbf{separazione degli insiemi},
con selezione del checkpoint sulla validazione e mai sul test; \textbf{isolamento della
variabile}, cioè nei confronti fra alternative cambia un solo elemento; e \textbf{dichiarazione
delle lacune}, elencate esplicitamente anziché omesse.

Tutte le valutazioni sono offline, su GPU NVIDIA RTX~4050 Laptop da 6\,GB. Nessuna misura dipende
da servizi remoti, quote o chiavi API; l'LLM è eseguito in locale. I comandi esatti sono in
Appendice~\ref{app:comandi}.

\section{Metriche}
\emph{Accuracy} è riportata per completezza ma non è la metrica primaria: su classi sbilanciate
premia i classificatori che ignorano le classi rare. La metrica primaria per stance e
fact-checking è la \textbf{Macro-F1}, media \emph{non pesata} delle F1 per classe: non pesando le
classi per frequenza, attribuisce a una classe rara e difficile lo stesso peso di una frequente,
così che un modello che ignorasse NEUTRAL, o che non emettesse mai SUPPORTED, venga penalizzato in
modo visibile. Poiché entrambi questi comportamenti si sono effettivamente manifestati, la scelta
si è rivelata diagnosticamente decisiva; le F1 per classe sono sempre riportate accanto alla
media, perché la media da sola può nascondere profili molto diversi.

Per il ranking si usano NDCG@$k$ (Equazione~\ref{eq:ndcg}), MAP e Recall@$k$; per la diversità
l'entropia normalizzata (Equazione~\ref{eq:diversity}); per l'accordo fra annotatori il $\kappa$ di
Cohen~\cite{cohen}, nella variante a \textbf{pesi quadratici} $w_{ij} = 1 - ((i-j)/(K-1))^2$, che
è quella corretta su scala ordinale: la versione non pesata tratterebbe ``0 contro 1'' e ``0
contro 2'' come lo stesso errore.

\section{Protocollo: stance detection}\label{sec:proto-stance}
Tre livelli di difficoltà e rilevanza applicativa crescenti. \textbf{Test held-out}: Macro-F1 e F1
per classe sui 4\,255 esempi di test, confrontate con la baseline zero-shot NLI sugli stessi
esempi. \textbf{Cross-dataset}: le stesse metriche calcolate \emph{separatamente} sui due
sottoinsiemi di provenienza, per verificare che la prestazione complessiva non sia trainata da un
solo dataset e che il modello generalizzi a target non visti. \textbf{Dominio applicativo}:
Macro-F1 sulle etichette gold dei corpora social, la misura più severa perché valuta il modello
fuori distribuzione, ed è anche l'unica condotta sui dati effettivamente elaborati in esercizio.

Va dichiarato che \textbf{non} esiste una valutazione equivalente per l'ecosistema news: mancano
etichette gold su articoli. La conseguenza per l'interpretazione della diversità è nella
Sezione~\ref{sec:validita-interna}.

\section{Protocollo: fact-checking}\label{sec:proto-fc}
La valutazione su ClimateFEVER è condotta a \textbf{due granularità} che misurano cose diverse.
\emph{Per-evidenza}: il motore predice l'etichetta di ogni coppia (claim, evidenza) e la si
confronta con l'etichetta della singola evidenza; è la misura pulita del motore come
classificatore, senza aggregazione. \emph{Claim-level}: il verdetto è ottenuto aggregando su tutte
le evidenze della claim e confrontato con l'etichetta complessiva; è il verdetto end-to-end, più
difficile perché somma all'errore di classificazione quello di aggregazione.

Il confronto fra motori è condotto a livello di claim e \textbf{a parità di evidenze}: a entrambi
vengono passate esattamente le stesse evidenze annotate del benchmark, così che l'unica differenza
sia il meccanismo che emette il verdetto. È il punto metodologico centrale: senza di esso una
eventuale superiorità dell'LLM sarebbe attribuibile in modo indistinguibile al motore o alla
catena di recupero. È inoltre prevista un'\textbf{ablazione} sulla dimensione del modello NLI
(\emph{base} contro \emph{large}), per verificare se il costo superiore sia giustificato.

\section{Protocollo: Stance Diversity}\label{sec:proto-div}
Per ciascun tema si eseguono le query canoniche su tutti gli ecosistemi disponibili, si registra la
diversità del pool per ogni coppia (query, ecosistema) e si riporta la media per ecosistema. La
procedura è identica per i due ecosistemi, il che elimina per costruzione i confondimenti di
trattamento. Resta il confondimento di campionamento della Sezione~\ref{sec:corpus-social}, che il
protocollo attuale \emph{non} controlla; il protocollo esteso prevede di riportare accanto a
ciascuna diversità anche quella a priori del corpus e il \emph{diversity lift}.

\section{Protocollo: retrieval}\label{sec:proto-retrieval}
\paragraph{Topic e pooling.} Le 11 query canoniche. Per ciascuna il pool è l'\textbf{unione} dei
primi 20 risultati di cinque run: BM25 puro, denso puro, RRF, fusione con $\alpha$ di default, e
l'insieme dei candidati riordinati dal reranker. I documenti mai entrati nel pool sono considerati
non rilevanti, secondo la convenzione di Cranfield.

L'unione multi-run non è un formalismo. Pescando dal solo run di fusione, i documenti che
esclusivamente BM25 o il denso propongono resterebbero non giudicati e conteggiati come non
rilevanti, sottostimando sistematicamente proprio le baseline che l'ablazione intende confrontare
--- è il \emph{pool bias}, e su un'ablazione di questo tipo ne invaliderebbe le conclusioni. La
misura sui dati reali quantifica il fenomeno: dei \textbf{380} giudizi del pool, \textbf{143 --- il
$38\%$ --- provengono da un solo run} (Tabella~\ref{tab:pool}).

\begin{table}[t]
\centering
\small
\begin{tabular}{lccc}
\toprule
\textbf{Tema} & \textbf{Query} & \textbf{Coppie nel pool} & \textbf{Da un solo run} \\
\midrule
Nucleare & 4 & 137 & 50 \\
Clima & 3 & 109 & 48 \\
Vaccini & 4 & 134 & 45 \\
\midrule
\textbf{Totale} & \textbf{11} & \textbf{380} & \textbf{143 (38\%)} \\
\bottomrule
\end{tabular}
\caption{Composizione del pool, unione di cinque run a profondità 20. L'ultima colonna quantifica
ciò che un pooling dal solo run di fusione avrebbe perso.}
\label{tab:pool}
\end{table}

\paragraph{Scala e annotazione.} Rilevanza graduata su tre livelli: $0$ non rilevante, $1$
marginalmente rilevante (tratta il tema ma non risponde alla query), $2$ pienamente rilevante. La
scala graduata è necessaria per l'NDCG e più informativa di quella binaria su un dominio in cui la
pertinenza è raramente netta. Ogni coppia è giudicata \textbf{due volte}, da due annotatori umani
che compilano fogli identici e separati: due colonne nello stesso file si contaminerebbero a
vicenda e l'accordo risulterebbe gonfiato. L'ordine delle righe è randomizzato entro ciascuna query e le colonne
di rango e punteggio non compaiono nel foglio: un annotatore che vede ``rank 1'' è primato a
giudicarlo rilevante, e il \emph{position bias} falserebbe proprio la metrica che si sta
costruendo. I criteri sono formalizzati in linee guida scritte
(Appendice~\ref{app:qrels}); la regola più rilevante --- e la più facile da violare su un corpus di
temi divisivi --- è la separazione fra \emph{rilevanza} e \emph{posizione}: un documento che nega
l'origine antropica del riscaldamento è pienamente rilevante per la query ``\emph{is climate change
caused by humans}''. Confondere i due assi significherebbe misurare il motore sulla sua capacità di
occultare le opinioni scomode, cioè l'opposto dell'obiettivo del sistema.

\paragraph{Sistemi confrontati.} Sei: BM25 puro; denso puro; RRF; fusione con $\alpha$ di default;
fusione con $\alpha$ tarato; fusione seguita da reranking. Il confronto fra le prime tre e la
quarta quantifica il guadagno della fusione pesata; fra la quarta e la quinta il guadagno della
taratura; fra la quarta e la sesta il guadagno del secondo stadio. La taratura di $\alpha$ avviene
in \textbf{leave-one-topic-out}: per ogni tema, $\alpha$ è scelto sulle query degli altri due e
valutato sulle proprie. Con 11 query, scegliere $\alpha$ sullo stesso insieme su cui si riporta il
risultato sarebbe sovradattamento garantito. Si riporta infine la \textbf{copertura dei giudizi},
cioè la frazione del top-$k$ effettivamente giudicata: se scende sensibilmente sotto 1 il confronto
è viziato dal pool e va dichiarato.

\section{Ciò che non viene valutato}
La \textbf{summarization} non è valutata quantitativamente: mancano riassunti di riferimento di
dominio e in loro assenza ROUGE e BERTScore non sono calcolabili. La \textbf{stance detection
sull'ecosistema news} non è valutata, per assenza di gold annotato su articoli. Il
\textbf{fact-checking end-to-end sui documenti reali} non è valutato: la valutazione usa le
evidenze fornite dal benchmark, quindi la qualità del recupero dalla KB propria del sistema non è
quantificata. Nessuna misura è accompagnata da intervalli di confidenza o test di significatività
(\S\ref{sec:validita-conclusione}).

% =================================================================
\chapter{Risultati}\label{cap:risultati}

\section{Stance detection}\label{sec:ris-stance}
Sul test held-out il modello fine-tunato raggiunge \textbf{Macro-F1 $0{,}771$}, contro $0{,}390$
della baseline zero-shot sugli stessi esempi (Tabella~\ref{tab:stance}).

\begin{table}[t]
\centering
\small
\begin{tabular}{lcccc}
\toprule
\textbf{Test held-out} ($n = 4255$) & \textbf{Macro-F1} & \textbf{FAVOR} & \textbf{AGAINST} & \textbf{NEUTRAL} \\
\midrule
Fine-tuned (DeBERTa-v3-base) & \textbf{0,771} & 0,715 & 0,740 & 0,858 \\
Baseline zero-shot NLI       & 0,390 & --- & --- & 0,125 \\
\bottomrule
\end{tabular}
\caption{Stance detection: Macro-F1 complessiva e F1 per classe.}
\label{tab:stance}
\end{table}

Il guadagno è distribuito su tutte le classi ma particolarmente marcato su \textbf{NEUTRAL}
($0{,}858$ contro $0{,}125$). Il fallimento della baseline lì è comprensibile: l'ipotesi
``\emph{This text is neutral about X}'' non è una buona ipotesi NLI, perché la neutralità è
l'assenza di una relazione piuttosto che una relazione, e un modello di entailment non ha modo
naturale di esprimerla. La conseguenza per il sistema è diretta: senza riconoscimento affidabile
di NEUTRAL la Stance Diversity sarebbe sistematicamente sovrastimata, perché i documenti neutri
verrebbero distribuiti fra FAVOR e AGAINST producendo distribuzioni artificialmente bipolari. La
capacità di riconoscere l'\emph{assenza} di posizione è quindi un prerequisito della metrica, non
un raffinamento.

\paragraph{Generalizzazione cross-dataset.} Macro-F1 $0{,}768$ su VAST e $0{,}741$ su SemEval. La
sostanziale parità è positiva per due ragioni: la prestazione complessiva non è trainata dal solo
dataset maggioritario, e --- più rilevante --- VAST contiene 763 target distinti nel test, in larga
parte non visti, per cui quel valore è una misura di generalizzazione \emph{cross-target} e non di
memorizzazione. È il dato che giustifica a posteriori la rinuncia a tecniche di adattamento di
dominio: la generalizzazione è stata ottenuta a monte, tramite la varietà dei target in
addestramento, anziché a valle tramite l'architettura.

\paragraph{Dominio applicativo.}\label{sec:ris-stance-social}
Sul gold dei corpora social la Macro-F1 scende a $\sim\!0{,}48$ (clima $0{,}478$, vaccini
$0{,}485$). Il calo, di circa 29 punti, \textbf{non è uniforme}: è quasi interamente imputabile a
NEUTRAL (F1 fra $0{,}07$ e $0{,}23$), mentre FAVOR e AGAINST restano fra $0{,}6$ e $0{,}7$. Il
profilo è quello di un modello che continua a distinguere bene le posizioni schierate e perde quasi
completamente la capacità di riconoscere l'assenza di posizione.

L'interpretazione proposta è che il ``neutro'' sui social sia un \textbf{costrutto diverso}. Nei
dataset di addestramento la neutralità è tipicamente un testo informativo che riporta fatti; nei
tweet, molti contenuti etichettati come neutri sono ironici, allusivi, frammentari o riferiti a un
contesto conversazionale assente. Non si tratta soltanto di uno spostamento di dominio linguistico
ma dello spostamento della \emph{definizione della classe}: un miglioramento richiederebbe una
ri-annotazione con linee guida specifiche, non semplicemente più dati.

\section{Fact-checking su ClimateFEVER}

\paragraph{Motore NLI e ablazione sulla dimensione.}\label{sec:ris-ablazione}
La misura per-evidenza dà Accuracy $0{,}606$ e Macro-F1 $0{,}449$ con il modello \emph{base}.
L'analisi per classe individua il collo di bottiglia: \textbf{SUPPORTED} ha F1 $\approx 0{,}20$,
contro $0{,}73$ di NEI e $0{,}42$ di REFUTED. Il profilo è quello di un modello
\emph{conservativo}: bravo a riconoscere l'assenza di evidenza, discreto nel confutare, debole nel
confermare --- e la debolezza persiste nella misura per-evidenza pura, dove le soglie del sistema
non intervengono, quindi è una proprietà del modello e non della politica di decisione.

La variante \emph{large} dà $0{,}620$ / $0{,}458$: $+1{,}4$ punti di accuracy a costo circa triplo.
Si adotta il modello \emph{base}. La conclusione metodologica è però più interessante di quella
operativa: che triplicare la dimensione produca meno di un punto e mezzo suggerisce che il limite
non risieda nella capacità del modello ma nella \emph{formulazione del compito} --- ipotesi che
l'esperimento successivo conferma.

\paragraph{NLI contro LLM a parità di evidenze.}\label{sec:ris-fc-confronto}
Sostituendo l'NLI con il ragionamento dell'LLM locale a livello di claim ($n = 1381$), i risultati
cambiano radicalmente (Tabella~\ref{tab:fc}): Macro-F1 da $0{,}350$ a $0{,}640$ ($+83\%$),
Accuracy da $0{,}393$ a $0{,}667$. Il salto decisivo è su SUPPORTED, da $0{,}065$ a $0{,}753$: il
ragionamento colma esattamente il punto in cui l'NLI falliva.

\begin{table}[t]
\centering
\small
\begin{tabular}{lccccc}
\toprule
\textbf{ClimateFEVER, claim-level} ($n{=}1381$) & \textbf{Acc} & \textbf{Macro-F1} & SUP & REF & NEI \\
\midrule
NLI (DeBERTa-v3-base-mnli-fever-anli) & 0,393 & 0,350 & 0,065 & 0,479 & 0,507 \\
\textbf{LLM (qwen2.5:7b, grounding puro)} & \textbf{0,667} & \textbf{0,640} & \textbf{0,753} & \textbf{0,565} & \textbf{0,601} \\
\midrule
$\Delta$ relativo & $+70\%$ & $+83\%$ & $+1058\%$ & $+18\%$ & $+19\%$ \\
\bottomrule
\end{tabular}
\caption{Fact-checking: confronto fra i due motori di verdetto a parità di evidenze.}
\label{tab:fc}
\end{table}

Tre osservazioni. Il confronto è \textbf{controllato}: entrambi i motori ricevono le stesse
evidenze annotate, quindi non è un confronto fra sistemi ma fra motori a parità di tutto il resto.
Il \textbf{grounding esclude la spiegazione più ovvia}, cioè che l'LLM risponda a memoria su claim
climatiche già viste; va detto con onestà che l'istruzione \emph{riduce} ma non elimina
formalmente il rischio, e verificarne l'assenza richiederebbe un esperimento di controllo con
evidenze contraffatte, non condotto qui. E il \textbf{profilo si ribalta}: l'NLI era pessimo nel
confermare e relativamente migliore nel confutare, l'LLM diventa il più forte nel confermare e
trova il suo punto debole in REFUTED. Non è quindi un modello uniformemente migliore ma due profili
di errore \emph{complementari}: riconoscere una conferma richiede un'inferenza a più passi ---
l'evidenza raramente ripete la claim, la implica attraverso passaggi impliciti --- mentre
riconoscere una contraddizione è in molti casi una questione di incompatibilità locale, più alla
portata di un classificatore addestrato su coppie di frasi. La complementarità suggerisce che
un'aggregazione dei due motori potrebbe superare entrambi.

Il guadagno ha un costo in latenza: circa $4$ secondi per claim contro decine di millisecondi, due
ordini di grandezza. Il trade-off è esplicito: l'NLI resta appropriato per l'uso interattivo, l'LLM
per l'elaborazione differita --- che è la strategia adottata nella dimostrazione, dove i verdetti di
entrambi i motori sono precalcolati sulle query canoniche, così che il confronto sia ispezionabile
dall'utente.

\paragraph{Una conferma collaterale sui documenti propri.} Precalcolare entrambi i motori sulle
stesse 11 query produce $160$ verdetti per motore, sugli stessi documenti e con le stesse evidenze.
Non sono una misura di accuratezza --- su questi documenti non esiste gold --- ma la
\emph{distribuzione} dei verdetti riproduce l'inversione di profilo misurata su ClimateFEVER: l'NLI
emette $2\%$ di SUPPORTED e $15\%$ di REFUTED, l'LLM ribalta le proporzioni a $17\%$ e $7\%$. Il
comportamento osservato sul benchmark non è dunque un artefatto di ClimateFEVER: si ripresenta fuori
campione, su un corpus giornalistico e su una knowledge base diversa. Entrambi restano prudenti ---
oltre tre quarti dei verdetti sono NOT\_ENOUGH\_EVIDENCE --- il che è il comportamento atteso di un
cancello di pertinenza tarato per non produrre falsi REFUTED.

\section{Stance Diversity: media contro social}\label{sec:ris-diversity}

\begin{table}[t]
\centering
\small
\begin{tabular}{lccc}
\toprule
\textbf{Stance Diversity} & \textbf{News} & \textbf{Social} & $H$ \textbf{a priori del corpus social} \\
\midrule
Energia nucleare & 0,545 & --- & --- \\
Cambiamento climatico & 0,322 & \textbf{0,621} & 1,000 \\
Vaccini & 0,229 & \textbf{0,600} & 1,000 \\
\bottomrule
\end{tabular}
\caption{Diversità di stance comparata. L'ultima colonna riporta l'entropia a priori del corpus
social, uniforme per costruzione (\S\ref{sec:corpus-social}).}
\label{tab:div}
\end{table}

Su entrambi i temi con corpus social la diversità misurata è circa doppia rispetto ai media. La
lettura immediata sarebbe controintuitiva rispetto al senso comune che associa i social alla
polarizzazione. Va però sottoposta a due qualificazioni di segno opposto.

\paragraph{Il risultato solido.} I valori dell'ecosistema news sono interpretabili in senso
assoluto: il corpus è un campione naturale, e $0{,}229$ sui vaccini significa che sulle query
canoniche il materiale recuperabile è fortemente concentrato su una sola posizione. Questo è
coerente con la misura indipendente della Sezione~\ref{sec:against-news}, che rileva circa il $3\%$
di documenti AGAINST nell'intero corpus. Due misure ottenute per vie diverse --- una sul pool
recuperato, una sull'intero corpus --- concordano, e la loro concordanza rafforza la conclusione
che \emph{la cronaca giornalistica su questi temi, nel campione considerato, è marcatamente
omogenea}. È il risultato più difendibile del capitolo.

\paragraph{Il confronto, invece, è confondato.} Il \emph{rapporto} fra i due ecosistemi non è
interpretabile come proprietà degli ecosistemi, per la ragione esposta nella
Sezione~\ref{sec:corpus-social}: la differenza somma un effetto reale ignoto e un artefatto di
campionamento. L'analisi completa, con la direzione dell'artefatto, è nella
Sezione~\ref{sec:validita-interna}.

\paragraph{Un'osservazione collaterale.} Il corpus social, pur avendo entropia a priori $1{,}000$,
produce pool con diversità $0{,}62$ e $0{,}60$, sensibilmente sotto il massimo. Lo scarto ha due
spiegazioni non mutuamente esclusive: il retrieval per pertinenza concentra le prospettive --- che
sarebbe un risultato interessante di per sé, misurabile con il \emph{diversity lift} --- oppure
l'errore di classificazione su NEUTRAL distorce la distribuzione misurata. Con i dati attuali non
è possibile separare i due contributi, ed è precisamente ciò che il lift, affiancato a un'analisi
di sensibilità sulla confidenza, permetterebbe di fare. Il valore social per il nucleare non è
disponibile, non essendo stato individuato un dataset di tweet comparabile su quel tema.

\section{Retrieval}\label{sec:ris-retrieval}

\subsection{Accordo fra annotatori}
Le 380 coppie del pool sono state giudicate \textbf{due volte, da due annotatori umani} --- i due
autori --- che hanno lavorato in cieco l'uno rispetto all'altro sulle stesse linee guida scritte
(Appendice~\ref{app:qrels}). Nessuna delle due serie è stata consultata durante la compilazione
dell'altra: l'accordo misurato su annotazioni allineate strada facendo sarebbe gonfiato e non
significherebbe nulla.

\begin{table}[t]
\centering
\small
\begin{tabular}{lccccc}
\toprule
& & & \multicolumn{3}{c}{\textbf{$\kappa$ di Cohen}} \\
\cmidrule(lr){4-6}
\textbf{Tema} & \textbf{$n$} & Accordo & non pesato & lineare & \textbf{quadratico} \\
\midrule
Energia nucleare & 137 & 0,839 & 0,744 & 0,836 & \textbf{0,904} \\
Cambiamento climatico & 109 & 0,761 & 0,587 & 0,712 & \textbf{0,814} \\
Vaccini & 134 & 0,784 & 0,657 & 0,770 & \textbf{0,856} \\
\midrule
\textbf{Complessivo} & \textbf{380} & \textbf{0,797} & 0,675 & 0,784 & \textbf{0,867} \\
\bottomrule
\end{tabular}
\caption{Accordo fra i due annotatori umani sui giudizi di rilevanza. Il valore da riportare su una
scala ordinale è quello a pesi quadratici; le altre due varianti sono date per confrontabilità con
la letteratura, che cita più spesso la forma non pesata.}
\label{tab:kappa}
\end{table}

Il $\kappa$ quadratico complessivo è $0{,}867$ --- \emph{almost perfect} sulla scala di
interpretazione convenzionale --- e non scende sotto $0{,}81$ su nessuno dei tre temi; il non
pesato vale $0{,}675$, \emph{substantial}. Il divario fra le due varianti non è un artificio a
proprio vantaggio ma il contenuto informativo della scelta: significa che il disaccordo, dove c'è,
è quasi tutto fra gradi \emph{adiacenti}.

È la struttura del disaccordo, più della sua media, a dire se le linee guida reggono. Sulle 380
coppie i due giudizi divergono 77 volte: 31 sul confine fra $0$ e $1$, 44 su quello fra $1$ e $2$, e
soltanto \textbf{2 a due gradini di distanza} --- i casi in cui un giudice dice ``non pertinente'' e
l'altro ``pienamente pertinente'', nessuno dei quali sul nucleare. Non ci sono cioè coppie su cui i
due annotatori abbiano letto documenti diversi: il disaccordo riguarda \emph{quanto} un documento
serva alla query, non \emph{se} vi risponda. Il tema climatico resta il più difficile ($\kappa$
quadratico $0{,}814$ contro $0{,}904$ del nucleare), coerentemente con query astratte su un corpus
tematicamente saturo, dove il confine fra ``parla del tema'' e ``risponde alla domanda'' è il più
sottile.

Le due serie non sono però intercambiabili, e la differenza è sistematica: il secondo annotatore
assegna il grado $2$ a $153$ coppie contro $117$ del primo, e in $39$ casi promuove a
``pienamente rilevante'' un documento che il primo ha giudicato marginale (il passaggio inverso
avviene $5$ volte). È una soglia di generosità diversa, non un disaccordo sul merito, e va tenuta
presente leggendo l'ablazione: alza i valori assoluti di tutte le metriche calcolate sulla seconda
serie senza spostare il confronto fra sistemi.

\subsection{Ablazione dei sistemi}
La Tabella~\ref{tab:ablazione-retrieval} riporta l'ablazione, calcolata due volte sulle due serie
di giudizi.

\begin{table}[t]
\centering
\small
\begin{tabular}{clcccccc}
\toprule
& & \multicolumn{3}{c}{\textbf{Annotatore 1}} & \multicolumn{3}{c}{\textbf{Annotatore 2}} \\
\cmidrule(lr){3-5}\cmidrule(lr){6-8}
\textbf{\#} & \textbf{Sistema} & NDCG & MAP & R@20 & NDCG & MAP & R@20 \\
\midrule
1 & BM25 puro                  & 0,616 & 0,488 & 0,583 & 0,643 & 0,494 & 0,604 \\
2 & Denso puro (bge)           & 0,656 & 0,538 & 0,701 & 0,719 & 0,563 & 0,699 \\
3 & Fusione RRF                & 0,653 & 0,535 & 0,671 & 0,704 & 0,548 & 0,673 \\
4 & Fusione, $\alpha{=}0{,}45$ & 0,674 & 0,556 & \textbf{0,702} & 0,723 & 0,571 & \textbf{0,708} \\
5 & Fusione, $\alpha$ tarato   & 0,644 & 0,538 & 0,685 & 0,719 & 0,574 & 0,700 \\
6 & Fusione $+$ reranking      & \textbf{0,715} & \textbf{0,577} & \textbf{0,702} & \textbf{0,779} & \textbf{0,601} & \textbf{0,708} \\
\bottomrule
\end{tabular}
\caption{Ablazione del retrieval su 11 query, calcolata sulle due serie di giudizi indipendenti.
NDCG è a cutoff 10; i run sono troncati alla profondità 20 a cui il pool è giudicato, quindi il MAP
è un MAP@20.}
\label{tab:ablazione-retrieval}
\end{table}

\paragraph{Quanto l'ordinamento dipende da chi ha giudicato.} Cinque sistemi su sei si ordinano allo
stesso modo sulle due serie --- \emph{rerank} $>$ \emph{fusione} $>$ \emph{denso} $>$ \emph{RRF} $>$
\emph{BM25} --- e su quei cinque il $\tau_b$ di Kendall vale esattamente $1{,}0$. Sui sei scende a
$0{,}733$, e l'intero scarto è imputabile a un solo sistema: la \emph{fusione tarata}, quinta sulla
prima serie e terza sulla seconda. Vale la pena guardare di quanto: sulla seconda serie supera il
denso di $0{,}0006$ di NDCG, cioè si colloca in un intervallo che 11 query non hanno la risoluzione
per distinguere. Il $\tau_b$ registra quel sorpasso come un'inversione perché è una statistica
d'ordine e ignora le distanze; la lettura corretta non è che il risultato sia instabile, ma che
\emph{è instabile esattamente sul sistema che questa sezione conclude non dover essere usato}
(\S\ref{sec:ris-retrieval}, taratura di $\alpha$). Le conclusioni sostantive --- il reranking come
contributo maggiore, la fusione sopra entrambi i rami singoli e sopra RRF, BM25 ultimo --- sono
invarianti rispetto all'annotatore.

I valori assoluti sono invece sistematicamente più alti sulla seconda serie, di $3$--$6$ punti di
NDCG: quell'annotatore assegna più giudizi di piena rilevanza ($153$ contro $117$) e c'è quindi più
guadagno disponibile nel numeratore dell'NDCG. È la ragione per cui \emph{i valori assoluti di un
NDCG non sono confrontabili fra collezioni giudicate da persone diverse}, mentre i confronti interni
alla stessa collezione lo sono. Salvo indicazione contraria i numeri commentati nel seguito sono
quelli della prima serie.

\paragraph{Il reranking è il contributo più grande.} Il secondo stadio porta l'NDCG@10 da $0{,}674$
a $0{,}715$ ($+6{,}1\%$ relativo; $+7{,}8\%$ sulla seconda serie), a Recall invariato per costruzione
--- riordina, non recupera. È coerente con l'attesa teorica e la conferma su dati propri. Il
dettaglio per tema (Tabella~\ref{tab:ablazione-topic}) mostra però che il guadagno non è distribuito:
viene quasi tutto dai vaccini ($0{,}599 \to 0{,}722$), mentre sul nucleare, dove la fusione è già a
$0{,}822$, il riordino è sostanzialmente neutro. Il reranker serve dove il primo stadio lascia
margine, non ovunque.

\paragraph{La fusione batte i rami singoli, ma il margine sul denso è sottile.} $0{,}674$ contro
$0{,}656$ del denso e $0{,}616$ di BM25: $+2{,}9\%$ relativo sul denso e $+9{,}4\%$ su BM25. Il
ramo denso porta la maggior parte del segnale e il contributo lessicale è reale ma modesto. La
fusione pesata supera anche RRF ($0{,}653$), il che giustifica il parametro: combinare i punteggi
conserva informazione che combinare i ranghi perde. Il dettaglio per tema
(Tabella~\ref{tab:ablazione-topic}) mostra la ragione della complementarità: sul nucleare, dove
pesano i termini tecnici, BM25 batte il denso ($0{,}800$ contro $0{,}737$); sul clima, dove domina la
parafrasi, crolla a $0{,}405$ contro $0{,}576$. È esattamente la complementarità ipotizzata in
progettazione, e la fusione serve a non dover scegliere. L'inversione fra i due rami è però il
risultato più fragile del capitolo: sulla seconda serie di giudizi il denso resta avanti anche sul
nucleare ($0{,}831$ contro $0{,}770$), e quello che sopravvive a entrambe le annotazioni è la
\emph{variazione del divario} --- il vantaggio del denso passa da $-6$ punti sul nucleare a
$+17$ sul clima nella prima serie, da $+6$ a $+17$ nella seconda --- non il segno.

\begin{table}[t]
\centering
\small
\begin{tabular}{lccc}
\toprule
\textbf{NDCG@10 (Annotatore 1)} & \textbf{Nucleare} & \textbf{Clima} & \textbf{Vaccini} \\
\midrule
BM25 puro & 0,800 & 0,405 & 0,591 \\
Denso puro & 0,737 & 0,576 & 0,634 \\
Fusione pesata & \textbf{0,822} & \textbf{0,579} & 0,599 \\
Fusione $+$ reranking & 0,820 & 0,567 & \textbf{0,722} \\
\bottomrule
\end{tabular}
\caption{Dettaglio per tema: il divario fra ramo sparso e denso si ribalta fra nucleare e clima, e
il guadagno del reranking è concentrato sui vaccini.}
\label{tab:ablazione-topic}
\end{table}

\paragraph{Tarare $\alpha$ su 11 query peggiora le prestazioni.} È il risultato più istruttivo. Il
sistema con $\alpha$ tarato in leave-one-topic-out ($0{,}644$) è \emph{inferiore} a quello con il
default ($0{,}674$); sulla seconda serie di giudizi la perdita si riduce a mezzo punto percentuale
($0{,}719$ contro $0{,}723$) ma non cambia di segno. La ragione è visibile nella curva completa
(Tabella~\ref{tab:alpha-curve}): fra $0{,}1$ e $0{,}5$ la metrica è praticamente piatta, con
un'escursione di $1{,}6$ punti, e degrada solo oltre $0{,}6$. L'ottimo \emph{oracolo} della prima
serie è $\alpha = 0{,}4$ con $0{,}6768$, cioè $+0{,}4\%$ rispetto al default: una differenza priva
di significato pratico. E i due annotatori collocano l'ottimo in punti diversi ($0{,}4$ contro
$0{,}2$), che è la firma del rumore. Con due soli temi di addestramento per fold la grid search
seleziona quel rumore ($\alpha = 0{,}1$ sul fold del nucleare, su entrambe le serie) e generalizza
peggio di un valore fissato a priori.

\begin{table}[t]
\centering
\small
\begin{tabular}{lccccccccc}
\toprule
$\alpha$ & 0,1 & 0,2 & 0,3 & 0,4 & \textbf{0,45} & 0,5 & 0,6 & 0,8 & 0,9 \\
\midrule
NDCG@10 (Annot.\ 1) & 0,661 & 0,673 & 0,667 & \underline{0,677} & \textbf{0,674} & 0,669 & 0,660 & 0,640 & 0,627 \\
NDCG@10 (Annot.\ 2) & 0,720 & \underline{0,731} & 0,725 & 0,721 & \textbf{0,723} & 0,712 & 0,710 & 0,692 & 0,672 \\
\bottomrule
\end{tabular}
\caption{NDCG@10 in funzione del peso del ramo sparso. In grassetto il default adottato,
sottolineato l'ottimo oracolo di ciascuna serie. La metrica è insensibile ad $\alpha$
nell'intervallo utile.}
\label{tab:alpha-curve}
\end{table}

La conclusione è duplice. Il default $\alpha = 0{,}45$, scelto per motivi di dominio, si rivela
entro lo $0{,}4\%$ dall'ottimo misurabile: la scelta a priori era buona, e ora si può dirlo sui
dati anziché sulla letteratura. Ma l'esperimento mostra anche che \emph{su questa scala di dati la
taratura di $\alpha$ non è un'operazione sensata}: non c'è segnale sufficiente a distinguere i
valori nell'intervallo utile, e tarare introduce varianza senza ridurre distorsione.

\paragraph{Copertura dei giudizi.} È $1{,}000$ per cinque sistemi su sei, come atteso: il pool è
l'unione dei loro run. L'unica eccezione è la fusione tarata ($0{,}986$ sulla prima serie, $0{,}982$
sulla seconda), che a $\alpha$ diverso dal default pesca qualche documento fuori dal pool; quei
documenti contano come non rilevanti e la penalizzano lievemente. L'effetto è di entità inferiore
al divario osservato ($3{,}0$ punti di NDCG) e non ne ribalta il segno, ma va dichiarato.

\section{Quadro d'insieme}

\begin{table}[t]
\centering
\small
\begin{tabularx}{\textwidth}{p{3.3cm}p{2.9cm}p{2.5cm}X}
\toprule
\textbf{Componente} & \textbf{Benchmark} & \textbf{Risultato} & \textbf{Stato} \\
\midrule
Retrieval, 6 sistemi & qrels propri (380) & NDCG@10 0,715 & valutato \\
Accordo annotatori & 2 umani, 380 coppie & $\kappa$ quad. 0,867 & \emph{almost perfect} \\
Stance detection & VAST $\cup$ SemEval & Macro-F1 0,771 & valutato \\
Stance, cross-dataset & VAST / SemEval & 0,768 / 0,741 & generalizza \\
Stance, dominio social & gold social & $\sim$0,48 & calo analizzato \\
Fact-check NLI & ClimateFEVER & Macro-F1 0,350 & valutato \\
Fact-check LLM & ClimateFEVER & Macro-F1 0,640 & confronto controllato \\
Stance Diversity & query canoniche & 0,229--0,621 & misurato, confronto confondato \\
Summarization & --- & --- & non valutabile \\
\bottomrule
\end{tabularx}
\caption{Stato di valutazione dei componenti.}
\label{tab:quadro}
\end{table}

% =================================================================
\chapter{Discussione e limiti}\label{cap:discussione}

\section{Findings trasversali}
\paragraph{L'omogeneità informativa è misurabile.} L'entropia normalizzata sulle stance predette è
calcolabile su qualunque corpus testuale, senza dati relazionali, e restituisce un numero
interpretabile. Il valore non sta nel particolare numero ottenuto ma nel fatto che l'echo chamber
diventi una grandezza monitorabile, aprendo alla possibilità di confronti sistematici nel tempo o
fra ambienti. Le condizioni sotto cui il confronto è lecito sono però più stringenti di quanto
appaia, ed è la lezione principale del lavoro.

\paragraph{Ragionare batte classificare, ma non ovunque.} A parità di evidenze un LLM locale di
taglia contenuta supera nettamente un classificatore NLI specializzato, e lo supera proprio dove
questo fallisce. Ma il profilo si ribalta su REFUTED. La conclusione corretta non è ``l'LLM è
migliore'', bensì che i due approcci hanno errori \emph{complementari} e che la scelta dipende da
quale errore sia più costoso --- o, meglio, che una loro combinazione meriterebbe di essere
esplorata.

\paragraph{Il secondo stadio vale più di qualunque taratura.} Sull'ablazione del retrieval, il
reranking --- che nessun iperparametro controlla --- vale un guadagno dieci volte superiore a
qualunque scelta di $\alpha$, la cui taratura su questa scala di dati è addirittura
controproducente. È un promemoria sull'ordine delle priorità in un sistema di IR.

\paragraph{Il grounding ha un prezzo e va pagato consapevolmente.} Vincolare il modello alle sole
evidenze recuperate elimina la fonte principale di allucinazione, ma limita la copertura a ciò che
la knowledge base contiene: una claim recente riceve NOT\_ENOUGH\_EVIDENCE non perché sia infondata
ma perché il sistema non sa. È corretto per uno strumento di verifica --- l'astensione è
preferibile all'affermazione arbitraria --- ma va comunicato all'utente, perché un NEI non è un
giudizio sulla claim.

\section{Analisi degli errori}
\paragraph{Il NEUTRAL come costrutto instabile.} Il filo conduttore di gran parte degli errori è la
classe NEUTRAL: riconosciuta eccellentemente sui benchmark ($0{,}858$) e quasi per nulla sui tweet
($0{,}07$--$0{,}23$), mentre le classi schierate reggono in entrambi i regimi. Lo schema esclude la
spiegazione più semplice --- un degrado generico fuori distribuzione, che colpirebbe tutte le classi
--- e indica che a cambiare è la \emph{definizione} della classe. L'implicazione per il sistema è
che la Stance Diversity misurata sul corpus social poggia, per una classe su tre, su una
classificazione poco più che casuale; e poiché un classificatore che sbaglia in modo non sistematico
tende a produrre distribuzioni più uniformi di quelle reali, l'errore agisce plausibilmente nella
direzione di \emph{gonfiare} la diversità misurata sul social.

\paragraph{SUPPORTED e la natura dell'inferenza.} Il fallimento dell'NLI su SUPPORTED ($0{,}065$ a
livello di claim) è quantitativamente il risultato più estremo del lavoro, e l'ablazione base/large
mostra che non si risolve con più capacità. La spiegazione proposta riguarda la natura del compito:
nei dati reali l'evidenza raramente \emph{ripete} la claim, la implica attraverso passaggi non
esplicitati (dall'evidenza ``la Francia produce il $70\%$ della sua elettricità da fonte nucleare''
alla claim ``la Francia dipende in larga misura dal nucleare'' occorre un passaggio di inferenza
lessicale e quantitativa). Un classificatore addestrato su coppie di frasi apprende bene la
contraddizione locale e male l'implicazione a più passi.

\paragraph{Falsi REFUTED e politiche di decisione.} Durante lo sviluppo si è osservato che un
fact-checker basato su recupero semantico su KB tematica produce falsi REFUTED con frequenza
elevata: le evidenze riguardano il tema giusto ma non sempre la claim giusta, e l'NLI interpreta la
divergenza di argomento come contraddizione. L'introduzione del cancello di pertinenza e
l'innalzamento della soglia di REFUTED a $0{,}95$ hanno convertito quei casi in NEI. È un esempio di
come una scelta di \emph{politica di decisione} possa incidere sul comportamento osservabile quanto
la scelta del modello.

\section{Minacce alla validità}

\subsection{Validità interna}\label{sec:validita-interna}
\paragraph{Minaccia principale: confondimento fra ecosistema e campionamento.} Il corpus social è
campionato in modo stratificato sulle tre classi (entropia a priori $=1{,}000$), quello news in
modo naturale (circa il $3\%$ di AGAINST). La differenza osservata confonde quindi due effetti:
\begin{equation}
\underbrace{H_{\text{social}} - H_{\text{news}}}_{\text{differenza osservata}} =
\underbrace{\Delta_{\text{ecosistema}}}_{\text{effetto di interesse}} +
\underbrace{\Delta_{\text{campionamento}}}_{\text{artefatto}} ,
\end{equation}
senza che i dati attuali consentano di separare i due termini. Sulla \emph{direzione}
dell'artefatto si può però ragionare con sicurezza: poiché la stratificazione porta l'entropia a
priori del social al massimo possibile mentre il campionamento naturale del news restituisce
un'entropia bassa, $\Delta_{\text{campionamento}} > 0$. L'artefatto agisce nella stessa direzione
della conclusione riportata e ne \emph{sovrastima} l'entità; non è però possibile escludere che
$\Delta_{\text{ecosistema}}$ sia nullo o di segno opposto.

Le contromisure, in ordine crescente di costo ed efficacia, sono: \emph{riformulare la
conclusione}, riportando i valori news come misura assoluta e presentando il confronto come
condizionato alla costruzione dei corpora --- soluzione adottata in questa relazione; applicare il
\emph{diversity lift} (\S\ref{sec:lift}), che rapportando la diversità del pool a quella del
proprio corpus è invariante rispetto alla distribuzione a priori e rende i due ecosistemi
confrontabili al costo di una passata offline; oppure \emph{ricostruire il corpus social con
campionamento naturale}, mantenendo quello stratificato come esclusivo gold di valutazione dello
stance detector --- soluzione definitiva, che elimina il confondimento alla radice.

\paragraph{Minaccia secondaria: errore di misura asimmetrico.} La Stance Diversity è calcolata su
etichette predette, e la qualità del predittore è misurata sul social ($\sim\!0{,}48$) ma
\textbf{non sul news}, per assenza di gold su articoli. È quindi possibile che l'errore di misura
sia di entità diversa nei due ecosistemi, e poiché l'errore tende a uniformare le distribuzioni,
una differenza di accuratezza si tradurrebbe in una differenza spuria di diversità. Contromisure:
annotare un gold pilota sul news (dell'ordine di 150 documenti bastano per una stima grossolana) e
condurre un'analisi di sensibilità ricalcolando la diversità sui soli documenti classificati con
confidenza superiore a una soglia.

\subsection{Validità esterna}
I corpora sono \textbf{piccoli}: 1917 articoli e 1869 tweet. Sufficienti a dimostrare il
funzionamento del sistema e a produrre misure indicative, non a caratterizzare l'ecosistema
mediatico. I \textbf{temi} sono tre, tutti sull'asse scienza/politica: nulla autorizza a estendere
le conclusioni ad altri assi di divisività. La \textbf{lingua} è esclusivamente l'inglese e tutti i
modelli sono monolingui. La \textbf{finestra temporale} è quella delle sorgenti, il che rende i
corpora una fotografia di un periodo e non una serie storica --- proprio la dimensione che
renderebbe interessante il monitoraggio. Va inoltre osservato che i due dataset social provengono da
fonti diverse con schemi di annotazione originari diversi, il che aggiunge eterogeneità fra i due
temi social.

\subsection{Validità di costrutto}
Il costrutto centrale è l'identificazione di \textbf{echo chamber} con \textbf{bassa entropia delle
stance nel materiale recuperato}: operazionalizzazione ragionevole ma parziale, perché l'echo
chamber riguarda l'esposizione effettiva di una persona nel tempo, mentre qui si misura la
composizione di un insieme di documenti. Una copertura omogenea può inoltre riflettere un consenso
scientifico legittimo anziché una distorsione informativa --- distinzione che la metrica non può
fare, e che è particolarmente rilevante proprio sui temi trattati. Un secondo aspetto è
l'\emph{invarianza per permutazione}, per cui una copertura prevalentemente neutra e una
prevalentemente schierata ricevono lo stesso punteggio: è il motivo per cui l'interfaccia mostra
sempre la distribuzione accanto al punteggio. Un terzo è la \textbf{dipendenza dal target}
(\S\ref{sec:target}), al punto che un target mal scelto può gonfiare la diversità del $50\%$ senza
che nulla sia migliorato: la metrica non è quindi proprietà del solo corpus, ma della coppia
(corpus, target).

\subsection{Validità delle conclusioni statistiche}\label{sec:validita-conclusione}
Nessuna misura è accompagnata da intervalli di confidenza o test di significatività. Per le misure
su benchmark ampi (stance su 4\,255 esempi, fact-checking su 1\,381 claim) la questione è meno
pressante: le differenze osservate sono di entità tale da rendere implausibile che siano dovute al
caso. Per la Stance Diversity, invece, il problema è concreto: le medie sono calcolate su 3--4
query per tema, ciascuna su un pool di 30 documenti, e con numerosità di questo ordine una
differenza di qualche centesimo non è distinguibile dal rumore. Lo stesso vale per l'ablazione del
retrieval, che poggia su 11 query: sufficiente a stabilire un ordinamento --- risultato corroborato
dall'invarianza di cinque sistemi su sei fra due serie di giudizi indipendenti --- ma non a
distinguere differenze piccole, come mostra proprio il sistema che si scambia di posto fra le due
serie a fronte di $0{,}0006$ di NDCG. Le contromisure appropriate sono riportare la deviazione
standard fra query e ampliare l'insieme delle query.

\section{Limiti}\label{sec:limiti}
\begin{description}[leftmargin=1.1cm,style=nextline]
  \item[Confronto fra ecosistemi confondato] La differenza di diversità fra social e news è
        sovrastimata per costruzione dei corpora. I valori assoluti sul news restano validi.
  \item[Numero di query ridotto nel retrieval] Undici query bastano a ordinare i sistemi ma non a
        tarare $\alpha$ né a dare intervalli di confidenza. Il parametro $b$ di BM25 resta un
        default motivato dal dominio e non tarato; l'esito della taratura di $\alpha$ rende
        plausibile, ma non dimostrato, che lo stesso valga per $b$.
  \item[Qrels non riconciliati] Le due annotazioni indipendenti servono a misurare l'accordo e la
        stabilità dell'ordinamento, ma le divergenze non sono state adjudicate in un'unica serie
        consolidata: i valori riportati restano quelli del primo annotatore, con il secondo come
        replica. Con soli due giudici, inoltre, il $\kappa$ è esso stesso una stima rumorosa.
  \item[Il NEUTRAL come costrutto] La nozione di neutralità sui social differisce da quella dei
        benchmark; la Macro-F1 nel dominio applicativo scende a $\sim\!0{,}48$ quasi interamente per
        questa ragione.
  \item[Stance sul news non misurata] Manca un gold annotato su articoli giornalistici.
  \item[Copertura della knowledge base] Claim recenti o molto specifiche producono NEI: è il limite
        intrinseco di una KB statica.
  \item[Fact-checking end-to-end non valutato] La valutazione usa le evidenze del benchmark, quindi
        la qualità del recupero dalla KB propria non è quantificata.
  \item[Grounding non verificato sperimentalmente] L'assenza di ricorso alla conoscenza parametrica
        dell'LLM è imposta via prompt ma non verificata con un esperimento di controllo.
  \item[Summarization non valutata] In assenza di riassunti di riferimento la qualità non è
        quantificata: è una baseline ragionevole, non un componente validato.
  \item[Metadati di fonte non sfruttati] Orientamento, affidabilità e tier sono calcolati e
        persistiti su ogni documento ma non utilizzati a valle; la copertura del registro sul corpus
        è inoltre bassa.
\end{description}

\section{Considerazioni etiche}\label{sec:etica}
\paragraph{Il rischio della falsa equivalenza.} Presentare sempre tre colonne affiancate può
suggerire che le tre posizioni abbiano pari fondamento. Su temi con un solido consenso scientifico
questa è precisamente la distorsione nota come \emph{false balance}. Il sistema adotta due
contromisure parziali: la barra di diversità mostra le proporzioni \emph{reali} del pool, rendendo
visibile che una posizione è minoritaria anziché nasconderlo dietro colonne di uguale ampiezza; e il
fact-checking affianca a ciascun documento un verdetto sulle sue affermazioni. Sono mitigazioni, non
soluzioni: il problema è intrinseco all'interfaccia a colonne ed è dichiarato.

\paragraph{Nessun filtro automatico sulle fonti.} Il registro di credibilità è usato per descrivere
e, in prospettiva, per segnalare, mai per escludere automaticamente documenti. Un sistema che
nascondesse silenziosamente determinate fonti eserciterebbe una curatela non trasparente, in
conflitto con l'obiettivo di esporre alla pluralità. Le fonti di tier 3 sono mostrate ma non sono
\emph{mai} usate come evidenza nel fact-checking: la distinzione fra ``ciò che si mostra'' e ``ciò
su cui si fonda un verdetto'' è la linea che il sistema traccia.

\paragraph{Reificazione delle posizioni e dati personali.} Ridurre un testo argomentato a una di tre
etichette è una semplificazione: il sistema conserva e mostra la confidenza, ma la struttura a tre
colonne impone comunque una categorizzazione netta, e la classificazione va intesa come strumento
di navigazione, non come giudizio sul testo. I corpora social contengono testi pubblici di utenti
reali: il sistema non raccoglie né mostra identificativi d'autore, non ricostruisce profili e non
collega testi allo stesso autore; i dataset impiegati sono pubblici e già anonimizzati alla fonte.

% =================================================================
\chapter{Conclusioni}\label{cap:conclusioni}

\section{Che cosa è stato ottenuto}
\ES{} realizza un approccio operativo e riproducibile alla misura dell'echo chamber. Dato un tema
divisivo, il sistema recupera i documenti pertinenti da due ecosistemi con una pipeline identica,
ne stima la posizione argomentativa, li bilancia per prospettiva, ne verifica le affermazioni
contro una base di conoscenza locale e li sintetizza per punto di vista, producendo una misura
scalare della diversità di opinione. Ripercorrendo le domande della Sezione~\ref{sec:rq}:

\begin{description}[leftmargin=1.3cm,style=nextline]
  \item[RQ1 (misurabilità)] \emph{Risposta affermativa}. La Stance Diversity è calcolabile su
        qualunque corpus testuale ed è stata analizzata nelle sue proprietà e nei suoi casi limite.
        Il lavoro ha inoltre identificato la condizione sotto cui il confronto fra due valori è
        lecito --- comparabilità delle procedure di campionamento --- e proposto una metrica
        derivata che ne è invariante.
  \item[RQ2 (confronto fra ecosistemi)] \emph{Risposta parziale e condizionata}. Il risultato solido
        è che l'ecosistema news dei corpora costruiti è marcatamente omogeneo in senso assoluto
        (diversità $0{,}23$--$0{,}55$; circa il $3\%$ di documenti AGAINST), con due misure
        indipendenti che concordano. Il confronto con il social è confondato dalla differenza di
        campionamento e non è interpretabile come proprietà degli ecosistemi finché non si applica la
        contromisura.
  \item[RQ3 (motore di verdetto)] \emph{Risposta affermativa e quantificata}. A parità di evidenze,
        su ClimateFEVER l'LLM locale porta la Macro-F1 da $0{,}350$ a $0{,}640$, con SUPPORTED da
        $0{,}065$ a $0{,}753$. Il profilo di errore dei due motori si rivela complementare: è il
        risultato più interessante e meno anticipato del lavoro.
  \item[RQ4 (fattibilità)] \emph{Risposta affermativa}. L'intera pipeline, che include sei modelli
        neurali distinti fra cui un LLM da 7 miliardi di parametri, è eseguibile su una GPU
        portatile da 6\,GB, interamente offline, con risultati riproducibili.
\end{description}

A queste si aggiunge un risultato non previsto fra le domande iniziali, emerso dall'ablazione del
retrieval: sul regime di dati considerato \textbf{la taratura di $\alpha$ non è un'operazione
sensata}, perché la metrica è piatta nell'intervallo utile del parametro e la procedura corretta ---
la validazione leave-one-topic-out --- peggiora le prestazioni selezionando rumore.

\section{Due lezioni trasversali}
La prima riguarda la \textbf{fragilità delle metriche derivate}. La Stance Diversity è elegante e
interpretabile, ma è calcolata su etichette predette da un classificatore e su un campione
costruito da una procedura. Il lavoro ha incontrato entrambe le dipendenze nella loro forma
concreta: un target mal scelto che gonfiava la diversità del $50\%$, e un campionamento stratificato
che ne determina in larga parte il valore su un intero ecosistema. Una metrica derivata va sempre
accompagnata dalla misura della qualità dei suoi ingredienti.

La seconda riguarda il \textbf{rapporto fra specializzazione e capacità generale}. Il confronto
NLI/LLM mostra che un modello specializzato e addestrato esattamente sul compito può essere superato
da un generalista di taglia contenuta, e superato proprio dove il primo fallisce; ma mostra anche
che la superiorità non è uniforme, e che là dove il compito si riduce al riconoscimento di
un'incompatibilità locale la specializzazione tiene. La lettura utile non è gerarchica ma
diagnostica: i due approcci sbagliano in modo diverso, ed è la struttura dell'errore --- non la
media --- a indicare quale usare.

\section{Sviluppi futuri}\label{sec:futuri}
\paragraph{Consolidamento.} \emph{Diversity lift}: calcolare $H(p^{\text{corpus}})$ per entrambi gli
ecosistemi e riportare il lift accanto alla diversità assoluta; è la contromisura al confondimento
principale e richiede una sola passata offline. \emph{Gold pilota sul news}, per stimare la Macro-F1
dello stance detector e valutare la simmetria dell'errore di misura fra ecosistemi. \emph{Più query}
(40--50), per intervalli di confidenza e test di significatività, e l'\emph{adjudicazione} delle 77
divergenze fra i due annotatori in un'unica serie consolidata, che è il passo che manca perché i
qrels smettano di essere il giudizio di una persona con una replica accanto.

\paragraph{Estensioni del sistema.} Sfruttare i \emph{metadati di fonte} già calcolati, mostrandoli
all'utente e affiancando alla Stance Diversity una \emph{Source Diversity} sull'orientamento delle
fonti del pool. \emph{Parallelizzare il ramo di fact-checking}, che è I/O-bound e non appartiene al
cammino critico. \emph{Persistere l'indice della knowledge base}, i cui embedding sono attualmente
ricalcolati a ogni avvio. Affiancare un \emph{RAG online} alla KB statica, mantenendo l'LLM locale
come motore di ragionamento, per verificare anche claim recenti.

\paragraph{Direzioni di ricerca.} \emph{Combinare i due motori di verdetto}, poiché i profili di
errore sono complementari: l'NLI come filtro rapido e l'LLM sui casi incerti, oppure una votazione
pesata per classe. Un \emph{confronto multi-LLM} fra modelli locali $\le$10B sul medesimo
protocollo. La \emph{verifica sperimentale del grounding} con evidenze contraffatte. La
\emph{ri-annotazione del NEUTRAL} con linee guida specifiche per il dominio social. E la
\emph{dimensione temporale}: corpora su più finestre per lo stesso tema, che trasformerebbero la
misura di diversità da fotografia in serie storica --- la forma in cui il monitoraggio dell'echo
chamber diventa davvero informativo.

% =================================================================
\appendix

\chapter{Iperparametri}\label{app:iper}
\begin{table}[h]
\centering
\small
\begin{tabular}{lll}
\toprule
\textbf{Sezione} & \textbf{Parametro} & \textbf{Valore} \\
\midrule
Corpus & chunk / overlap & 512 / 128 token (25\%) \\
\midrule
BM25 & $k_1$ / $b$ & 1,2 / 0,60 \\
\midrule
Dense & modello & \texttt{BAAI/bge-base-en-v1.5} (768-dim) \\
 & indice & FAISS \texttt{IndexFlatIP}; max-pooling sui chunk \\
\midrule
Fusione & $\alpha$ / $\beta$ & 0,45 / 0,55, normalizzazione min-max \\
 & pool & 50 \; (baseline RRF con $k_{\text{rrf}} = 60$) \\
\midrule
Reranker & modello / top-$n$ & \texttt{BAAI/bge-reranker-base} / 30 \\
\midrule
Stance & modello & DeBERTa-v3-base fine-tunato (v.\ Tab.~\ref{tab:iperparam-training}) \\
 & batch inferenza / troncamento & 16 / 400 parole \\
\midrule
Colonne & documenti per colonna & 5 \\
\midrule
Fact-check & KB & 1589 articoli Wikipedia (introduzioni) \\
 & cancello di pertinenza & coseno $\ge 0{,}62$; 4--5 evidenze per claim \\
 & NLI & \texttt{DeBERTa-v3-base-mnli-fever-anli} \\
 & soglie & SUPPORTED $\ge 0{,}75$; REFUTED $\ge 0{,}95$ \\
 & LLM & \texttt{qwen2.5:7b} (Ollama), temperatura 0, timeout 180\,s \\
\midrule
Summarization & modello & \texttt{google/flan-t5-base} \\
 & frasi estratte / beam & 10 / 4 \; (\texttt{do\_sample=False}) \\
 & max\_new\_tokens / no\_repeat & 200 / 3 \\
\bottomrule
\end{tabular}
\caption{Iperparametri del sistema. Nel percorso interattivo beam e token sono ridotti (2 / 120).}
\end{table}

\chapter{Prompt del motore LLM}\label{app:prompt}
Prompt di sistema effettivamente usato. Si notino l'istruzione di grounding puro, la definizione
esplicita delle tre classi e il vincolo di formato che rende la risposta parsabile.

\begin{lstlisting}
You are a rigorous fact-checking assistant. You receive a CLAIM and a
numbered list of EVIDENCE passages. Using ONLY the evidence (never prior
knowledge), decide:
- SUPPORTED: the evidence confirms/entails the claim;
- REFUTED: the evidence contradicts the claim;
- NOT_ENOUGH_EVIDENCE: the evidence is insufficient, unrelated, or ambiguous.
Reply ONLY with JSON: {"verdict": "SUPPORTED|REFUTED|NOT_ENOUGH_EVIDENCE",
"reasoning": "<one concise sentence citing the evidence used>"}.
\end{lstlisting}

Il messaggio utente ha la forma \texttt{CLAIM: ...} seguita da \texttt{EVIDENCE:} e dall'elenco
numerato delle evidenze recuperate. Il parser normalizza le varianti che i modelli producono
spontaneamente (\texttt{SUPPORTS}, \texttt{TRUE}, \texttt{ENTAILMENT} $\to$ SUPPORTED;
\texttt{REFUTES}, \texttt{FALSE}, \texttt{CONTRADICTION} $\to$ REFUTED; \texttt{NEI},
\texttt{INSUFFICIENT}, \texttt{UNKNOWN} $\to$ NOT\_ENOUGH\_EVIDENCE) e ripiega su
NOT\_ENOUGH\_EVIDENCE in assenza di segnale o in caso di errore di rete.

\chapter{Linee guida di annotazione dei qrels}\label{app:qrels}
Si annota la \textbf{rilevanza del documento rispetto alla query}: non la sua stance, non la sua
veridicità, non la qualità della fonte. La regola che governa tutte le altre è che
\textbf{rilevante non significa condivisibile né vero}: un articolo che nega l'origine antropica
del riscaldamento, se risponde alla query ``\emph{is climate change caused by humans}'', è
pienamente rilevante.

Il punto critico della scala è il confine fra $0$ e $1$. Definirlo per intenzione (``parla del tema
ma non ha rapporto con la query'' contro ``tratta il tema ma non risponde'') lo lascia
operativamente ambiguo, e l'ambiguità morde esattamente dove la query è astratta e il corpus
tematicamente saturo --- cioè sulle query climatiche, che restano infatti il tema di minore accordo
(Tabella~\ref{tab:kappa}). La specifica è perciò formulata come \emph{test verificabile sul testo},
non come descrizione di un'intenzione.

\begin{center}
\small
\begin{tabularx}{\textwidth}{cX}
\toprule
\textbf{Voto} & \textbf{Test operativo} \\
\midrule
\textbf{2} & Si sa \emph{indicare la frase} del documento che costituisce la risposta alla query. \\
\textbf{1} & Non c'è una frase-risposta, ma il documento dà contesto utile o tocca la domanda di
lato: lo si aprirebbe come terzo o quarto risultato ricavandone qualcosa. \\
\textbf{0} & Il legame con la query è \emph{solo tematico}: il documento parla del tema ma non della
domanda. Include il fuori tema puro. \\
\bottomrule
\end{tabularx}
\end{center}

\noindent\textbf{Corollari, dedotti da casi reali del corpus.}
\begin{itemize}[leftmargin=*]
  \item Se la query chiede un \textbf{costo}, serve un costo: una cifra, una stima o
        un'affermazione esplicita su chi paga e quanto. Una disputa su trasferimenti di fondi fra
        enti pubblici non è il costo di una politica; e il costo \emph{dei danni} non è il costo
        \emph{della politica}, hanno segno opposto.
  \item Se la query chiede una \textbf{relazione} fra due cose
        (clima\,$\leftrightarrow$\,eventi estremi, nucleare\,$\leftrightarrow$\,clima), il testo
        deve \emph{asserire} la relazione: non basta che entrambi i termini compaiano.
  \item Se la query chiede un \textbf{rischio} o un \textbf{effetto}, il documento deve parlare di
        quel rischio, non del tema in cui il rischio esiste. Il fallimento finanziario di un
        progetto non è un ``rischio dell'energia nucleare'' se il testo non lo inquadra così.
  \item Il \textbf{duplicato sindacato} prende lo stesso voto della sua copia: il corpus contiene
        molte riprese della stessa storia d'agenzia, che sono un'unica decisione.
  \item In dubbio fra $1$ e $0$ si scelga $0$; fra $2$ e $1$ si scelga $1$.
\end{itemize}

\noindent\textbf{Procedura.} Ogni annotatore compila il proprio foglio senza consultare l'altro
durante il lavoro: l'accordo misurato su annotazioni allineate strada facendo è gonfiato e non
significa nulla. Il $\kappa$ si calcola \emph{prima} di qualunque riconciliazione; solo dopo si
adjudicano le divergenze per produrre i qrels definitivi. Del protocollo è eseguito il primo passo
--- le due annotazioni cieche e il $\kappa$ (Tabella~\ref{tab:kappa}) ---; l'adjudicazione resta fra
gli sviluppi (\S\ref{sec:futuri}), e i risultati riportati usano la serie del primo annotatore.

\chapter{Riproducibilità}\label{app:comandi}
Il sistema è eseguibile a partire dagli indici pre-costruiti; la demo carica risultati precalcolati
per le query canoniche, così da essere istantanea e non richiedere GPU a runtime.

\begin{lstlisting}
# ambiente
uv sync --extra dev
uv run python -m spacy download en_core_web_sm

# costruzione dei dati (ogni script produce un MANIFEST.json con hash SHA-256)
uv run --extra data python scripts/build_corpora.py        # corpus news (CC-News)
uv run python scripts/build_social_corpus.py               # corpus social (tweet)
uv run python scripts/build_kb.py                          # KB Wikipedia
uv run python scripts/build_stance_splits.py               # split VAST + SemEval
uv run --extra data python scripts/index_corpus.py         # indici BM25 + FAISS
uv run python scripts/build_qrels_annotation.py            # fogli qrels da annotare

# addestramento e valutazione
uv run python scripts/train_stance.py                      # fine-tuning stance
uv run python scripts/eval_stance.py                       # benchmark stance
uv run python scripts/evaluate_pivot.py                    # stance social + diversity
uv run python scripts/eval_climatefever.py --engine nli    # fact-check, motore NLI
uv run python scripts/eval_climatefever.py --engine llm    # fact-check, motore LLM
uv run python scripts/eval_retrieval.py                    # ablazione retrieval
uv run python scripts/eval_agreement.py                    # kappa fra i due annotatori

# esecuzione e test
uv run python main.py "risks of nuclear energy" --topic nuclear_energy
uv run streamlit run dashboard.py
uv run pytest tests -q                                     # 174 test
\end{lstlisting}

% =================================================================
\begin{thebibliography}{99}
\bibitem{bge} S. Xiao, Z. Liu, P. Zhang, N. Muennighoff. \emph{C-Pack: Packed Resources for General
  Chinese Embeddings (BGE)}. 2023.
\bibitem{bm25} S. Robertson, H. Zaragoza. \emph{The Probabilistic Relevance Framework: BM25 and
  Beyond}. Foundations and Trends in Information Retrieval, 2009.
\bibitem{mmr} J. Carbonell, J. Goldstein. \emph{The Use of MMR, Diversity-Based Reranking for
  Reordering Documents and Producing Summaries}. Proc. SIGIR, 1998.
\bibitem{cinelli} M. Cinelli, G. De Francisci Morales, A. Galeazzi, W. Quattrociocchi, M. Starnini.
  \emph{The Echo Chamber Effect on Social Media}. PNAS, 2021.
\bibitem{claimbuster} N. Hassan, F. Arslan, C. Li, M. Tremayne. \emph{Toward Automated
  Fact-Checking: Detecting Check-worthy Factual Claims by ClaimBuster}. Proc. KDD, 2017.
\bibitem{climatefever} T. Diggelmann, J. Boyd-Graber, J. Bulian, M. Ciaramita, M. Leippold.
  \emph{CLIMATE-FEVER: A Dataset for Verification of Real-World Climate Claims}. 2020.
\bibitem{cohen} J. Cohen. \emph{A Coefficient of Agreement for Nominal Scales}. Educational and
  Psychological Measurement, 1960.
\bibitem{deberta} P. He, J. Gao, W. Chen. \emph{DeBERTaV3: Improving DeBERTa using ELECTRA-Style
  Pre-Training with Gradient-Disentangled Embedding Sharing}. Proc. ICLR, 2023.
\bibitem{dpr} V. Karpukhin et al. \emph{Dense Passage Retrieval for Open-Domain Question Answering}.
  Proc. EMNLP, 2020.
\bibitem{faiss} J. Johnson, M. Douze, H. Jégou. \emph{Billion-scale Similarity Search with GPUs}.
  IEEE Transactions on Big Data, 2019.
\bibitem{fever} J. Thorne, A. Vlachos, C. Christodoulopoulos, A. Mittal. \emph{FEVER: a Large-scale
  Dataset for Fact Extraction and VERification}. Proc. NAACL, 2018.
\bibitem{pariser} E. Pariser. \emph{The Filter Bubble: What the Internet Is Hiding from You}.
  Penguin Press, 2011.
\bibitem{rag} P. Lewis et al. \emph{Retrieval-Augmented Generation for Knowledge-Intensive NLP
  Tasks}. Proc. NeurIPS, 2020.
\bibitem{rrf} G. V. Cormack, C. L. A. Clarke, S. Buettcher. \emph{Reciprocal Rank Fusion Outperforms
  Condorcet and Individual Rank Learning Methods}. Proc. SIGIR, 2009.
\bibitem{sbert} N. Reimers, I. Gurevych. \emph{Sentence-BERT: Sentence Embeddings using Siamese
  BERT-Networks}. Proc. EMNLP, 2019.
\bibitem{semeval} S. Mohammad, S. Kiritchenko, P. Sobhani, X. Zhu, C. Cherry. \emph{SemEval-2016
  Task~6: Detecting Stance in Tweets}. Proc. SemEval, 2016.
\bibitem{shannon} C. E. Shannon. \emph{A Mathematical Theory of Communication}. Bell System
  Technical Journal, 1948.
\bibitem{vast} E. Allaway, K. McKeown. \emph{Zero-Shot Stance Detection: A Dataset and Model using
  Generalized Topic Representations (VAST)}. Proc. EMNLP, 2020.
\bibitem{yin} W. Yin, J. Hay, D. Roth. \emph{Benchmarking Zero-shot Text Classification: Datasets,
  Evaluation and Entailment Approach}. Proc. EMNLP, 2019.
\end{thebibliography}

\end{document}
